\documentclass{report}
\usepackage{amsmath, amssymb}
\usepackage{amsthm}
\usepackage{titlesec}
\usepackage[margin=1in]{geometry}
\usepackage{titling}
\usepackage{tabularx}
\usepackage{array}
\usepackage{pdflscape}
\usepackage{caption}
\usepackage{subcaption}
\usepackage{float}
\usepackage{graphicx}
\usepackage{cite}
\usepackage{rotating}
\usepackage{fancyhdr}
\usepackage{enumitem}
\usepackage{hyperref}
\usepackage{pifont}
\usepackage{array}
\usepackage[T1]{fontenc}
\usepackage{lmodern}
\usepackage{kpfonts}

% Define a custom page style for landscape pages
\fancypagestyle{landscapeStyle}{
  \fancyhf{}
  \cfoot{\thepage}
  \renewcommand{\headrulewidth}{0pt}
  \renewcommand{\footrulewidth}{0pt}
}

\newtheorem{theorem}{Theorem}
\newtheorem{definition}{Definition}
\newtheorem{lemma}{Lemma}
\newtheorem{corollary}{Corollary}
\newtheorem{proposition}{Proposition}

\newcolumntype{Y}{>{\raggedright\arraybackslash}X}
% Define custom numbering format
\titleformat{\section}{\normalfont\Large\bfseries}{\thesection}{1em}{}
\titleformat{\subsection}{\normalfont\large\bfseries}{\thesubsection}{1em}{}
\titleformat{\subsubsection}{\normalfont\normalsize\bfseries}{\thesubsubsection}{1em}{}

% Customize the numbering scheme
\renewcommand\thesection{\Roman{section}}
\renewcommand\thesubsection{\Alph{subsection}}
\renewcommand\thesubsubsection{\roman{subsubsection}}

% Define subtitle command
\newcommand{\subtitle}[1]{%
\posttitle{%
  \par\end{center}
  \begin{center}\large#1\end{center}
  \vskip0.5em}%
}

\pretitle{\begin{center}\LARGE\bfseries}
\posttitle{\par\end{center}\vskip 0.5em}
\preauthor{\begin{center}\large}
\postauthor{\end{center}}
\predate{\begin{center}\large}
\postdate{\end{center}}
% Redefine \cleardoublepage to \clearpage to avoid blank pages
\let\cleardoublepage\clearpage

% Redefine \part to use \clearpage instead of \cleardoublepage
\makeatletter
\renewcommand\part{\if@openright\clearpage\else\clearpage\fi
  \thispagestyle{plain}%
  \if@twocolumn
    \onecolumn
    \@tempswatrue
  \else
    \@tempswafalse
  \fi
  \null\vfil
  \secdef\@part\@spart}
\makeatother

% Define the universal header style
\fancypagestyle{universalStyle}{
  \fancyhf{}
  \fancyhead[L]{The Entropy Game}
  \fancyhead[C]{\leftmark}
  \fancyhead[R]{\thepage}
  \renewcommand{\headrulewidth}{0.4pt}
  \renewcommand{\footrulewidth}{0pt}
}

% Apply the universal style to all pages
\pagestyle{universalStyle}

\begin{document}

\begin{titlepage} % Start titlepage environment
    \centering
    \vspace*{\fill} % Vertically center
    \Huge\bfseries The Entropy Game\\
    \Large Physics is NP-Complete \& P=NP\\
    \vspace{1cm}
    \Large Essam Abadir\\
    \vspace{1cm}
    \large \today\\

    \vspace{3cm}
    \normalfont \emph{In memory of Gian-Carlo Rota, April 27, 1932 - April 18, 1999.}
    \vspace*{\fill} % Vertically center
    
\end{titlepage}

\tableofcontents
\clearpage % Ensure the table of contents ends cleanly

\part{Beforehand}
\markboth{Beforehand}{Beforehand} % Set the header for the part

\section*{60 Pages ... In 60 Seconds!}
\addcontentsline{toc}{section}{60 Pages ... In 60 Seconds!}
I've done enough AI ad targeting to know to include clickbait. For many, this will be the first and last thing you read, so here is what is absolutely essential to know without all the backup, careful math, or any restraint.

\subsection*{We're Going To Solve The Two Biggest Problems In Science: Quantum Gravity \& $P=NP$ ... And They Are The Same Problem!}
The two biggest problems in science are whether \(P=NP\) and how to reconcile Quantum Mechanics with General Relativity. The first problem is a math problem about whether we can solve almost \emph{all} really hard problems really quickly - it would be like jumping a century ahead in computing power. The second is a physics problem about how the universe works at the most fundamental level. No progress has been made on either one for over 50 years\cite{jaimungal2021wrong}\cite{jaimungal2021crisis}\cite{jaimungal2021gravitise}\cite{50yearsofPvsNP}\cite{Rota1985}\cite{Aaronson2013}\cite{aaronson2021odds}\cite{jaimungal2021greatest}. I will show that they are the same problem and that they can both be solved efficiently. If I'm being honest, all the combinatoric math I present is math I did over 30 years ago for problem sets in a class I took in 1993 by Professor Gian-Carlo Rota - it's his math.

\subsection*{Why Has No One Done This Before? They'll Say I'm Crazy, But I'm Not The One Saying 1+1=4!}
Physicists and theoretical mathematicians really love calculus and "fractions" (continuous numbers). The problem with believing in calculus as something real is that it leads to absurdities like "Schrödinger's Cat" being both dead and alive, the moon not existing when you aren't looking at it, parallel universes, time not existing, particles being in two places at once, and even quantum computer math where 1+1 can equal 4. The problem with believing in fractions is that there is no single example in nature of a fraction\cite{Gisin2018}. Every time we look closer at something we see that it is made up of whole number parts - the galaxy is made up of solar systems, the solar system has stars and planets, the planets have moons, the moons have rocks, the rocks have atoms, the atoms have electrons, protons, and neutrons, and the electrons have quarks. Physicists’ love of calculus is so great that they turned "Quantum Mechanics" into an oxymoron - quantum means "discrete amount" in Greek - but they decided to use continuous calculus for it instead! Of course, if you use the wrong math you get the wrong answers.

\subsection*{Quantum Means Discrete In Greek, Call it Discrete Mechanics}
The word "quantum" comes from the Latin word for "how much" and the Greek word for "discrete amount." Continuous is the antonym of discrete. The word "quantum" is used in physics to describe the smallest possible discrete unit of any physical property, such as energy or matter. The discovery of Quantum Mechanics was a revolution in physics because it showed that the universe is made up of discrete parts - we could as easily call "Quantum Mechanics" "Discrete Mechanics."

One would think that discrete mathematics should be used for Discrete Mechanics, but in fact, they use only calculus and continuous mathematics. 

\subsection*{Gravity + Calculus = Chaos = No Bueno}
The most important thing to know is that gravity is a non-linear system, and calculus is linear. The field of Chaos Theory was quite literally started with the gravitational Three-Body Problem. In the 1890s, Henri Poincaré showed that calculus cannot be used to solve gravity problems for three or more bodies (like sun, moon, and earth). So, we've known for over 130 years that gravity + calculus = No Bueno. Ironically, the field of Chaos Theory only uses calculus!?!

\subsection*{I'm On Team Einstein, Not Team Bohr!}
Einstein invented Discrete Mechanics, but his nemesis, Niels Bohr, insisted on using calculus, leading to the “Copenhagen interpretation,” where one thing can exist in two places at once—because, why not? In 1935, Einstein’s EPR paper called out this absurdity, but he was ignored. In the addendum, I give a math proof from another paper of mine that shows that Einstein was right and Bohr was wrong. How do I prove Bohr was wrong? The Bohr (and later the Bell Inequality) argument against Einstein is based on linear calculus and makes believe that chaos (non-linearity) doesn't exist in the quantum realm. I show that if gravity exists in the quantum realm, then chaos must be present and therefore any linear calculus-based argument against Einstein is flawed because ... gravity + calculus = chaos = no bueno.

\subsection*{Quantum Gravity, The Biggest Unsolved Physics Problem - Solved \checkmark}
So, to review. Quantum ("Discrete") Mechanics only uses calculus. If you had to guess whether it was possible to bring gravity into "Discrete Mechanics" using calculus, what would you now guess? For the past 100 years, Discrete Mechanics has clung to calculus while trying to add gravity into the theory … but, calculus + gravity = no Bueno. No surprise, it doesn’t work!

Combinatorics is often referred to as "balls into boxes" counting problems. Without Quantum Gravity (boxes) and discrete energy units (balls that go into the boxes) you have no counting problem to which you can apply Combinatorics. Once you introduce those things (i.e. you add gravity to Quantum Mechanics) the math is relatively easy - an AP Stats high school student could definitely follow it. A full detailed mathematical derivation of Quantum Gravity à la the Combinatorics I would have done in Rota's class is presented in the addendum. Physics' biggest problem \checkmark.

\subsection*{The Greatest Unsolved Problem in Mathematics: Does $P=NP$? - Solved \checkmark}
The question of whether \(P=NP\) is considered one of the most important open problems in mathematics and computer science. The "P" class represents problems that can be solved efficiently by a computer, in polynomial time. The "NP" class includes problems whose solutions can be quickly verified (also in polynomial time), but finding those solutions might be extremely difficult, perhaps even impossible, for a computer. Whether these two classes are the same (i.e., $P=NP$) would imply that any problem whose answer can be quickly checked can also be quickly solved. Most mathematicians and computer scientists suspect that \(P \neq NP\), but no one has been able to prove it either way, leaving it one of the Millennium Prize Problems with a million-dollar reward for a correct proof. The ramifications for computing and the world are so profound that if \(P=NP\) it would be comparable to the scientific revolution of the late 1800s when electricity was harnessed- many have even described it as the "end of science" in the sense that once everything can be efficiently solved, there are few problems left to explore.

\subsection*{Everything I Present Here Was Already Known in 1980}
From 1966 to 1999, professor Gian-Carlo Rota taught 18.313 which was innocuously called "Introduction to Probability Theory." The course was part of his influential work in establishing combinatorics as a serious mathematical discipline. His lectures and teaching style left a lasting impact on the field and on generations of students. I took the class in 1993 and what I realize only now is that it gave the first and only formal bridge between continuous concepts and discrete mathematics - most commentators still say this is impossible\cite{HossenfelderLorentz2015}. In fact, the vast majority of the math in here is mainly a "first publishing" of Chapter VIII of the unpublished manuscript from which he taught his class-I don't think that either I or any of my classmates knew that he was teaching us something truly unique. I can't say for certain when he established the formal bridge between these different areas of math but it dates back to at least 1979 because another student posted their earlier 1979 draft of the same text on the Internet Archive. Sadly, his passing in 1999 meant that this knowledge has since been hidden from the world.

Here's a recap of what 1980 was like in math, physics, and computer science:
Many if not most great physicists currently say that fundamental physics has made no real advances since the 1980s. In 1980, disagreeing with Bohr was career suicide. In 1980 P=NP was barely known and its implications were not well understood beyond theoretical computer science. The RSA public key encryption algorithm, which hinges on the difficulty of factoring large numbers, hadn't been widely published and its implications for cryptography and secure communications were also not widely understood. The internet as we know it existed only on research networks and the desktop computer was barely the equivalent of today's calculator. Even in 1993, complex Combinatorics made calculations of the type that were possible in Rota's 18.313 Combinatorics class unlikely to even be attempted. In short, the very math that could have transformed computing was being developed at the same time that computers were too weak to use it effectively. Instead, the prevailing winds of physics still favored continuous mathematics and calculus, and the digital computer as a scientific tool was not yet in wide use. This all combined to form a "perfect storm" of neglect, where a discrete and combinatorial approach to physics and computation was mostly ignored in favor of the prevailing ideas of continuous calculus.

He showed that if something exhibits "entropy-like" behavior (even those things physicists think of as continuous waves), it's actually a coded form of digital information that can be turned into an efficient computer program. Here's the math in a nutshell:

\begin{itemize}
    \item \textbf{Rota's Entropy Theorem: The Universal Translator}: Rota's theorem shows that all those "continuous" entropy distributions that physicists love? They are mathematically equivalent to a scaled version of digital Shannon Entropy. Think of it as a universal translator between the two worlds.
    \item \textbf{Shannon's Coding Theorem: The Fast Track}: When you have a Shannon Entropy problem (and Rota says you *always* do when looking at "entropy") Shannon's Coding Theorem guarantees that a fast computer program can be written to compute it, and it works even better as the problem gets more massive (like all real-world physical systems of particles, atoms, stars in a galaxy, etc.). It guarantees a solution that can be found in a time complexity of \(O(N \log N)\) - i.e. really fast, faster than the fastest theoretical "quantum computer" could ever do it.
    \item \textbf{The \(O(N \log N)\) Implies \(P=NP\): The Big Reveal}: Solving a math problem called $P=NP$ would transform the world, this is not hyperbole. Because of Shannon's Coding Theorem, a problem that can be solved in \(O(N \log N)\) time belongs to the "P" class, not the "NP" class. This implies that \(P=NP\). It is even more transformative if Physics is NP-Complete because it means every real-world problem is a physics problem and can be solved in \(O(N \log N)\) time.
\end{itemize}

So, what does all this mean? It means that the key to understanding any physics problem is to understand its discrete building blocks and the "code" that links them. Nature has been "coding" solutions all along.

This "lost" theorem is at the heart of everything we will be discussing, including the claim that \(P=NP\) and that physics is NP-Complete.

Now you want to read the rest of the paper, right?

\section*{Foreword: The Mathematician \& His Lost Theorem}
\addcontentsline{toc}{section}{Foreword: The Mathematician \& His Lost Theorem}

\begin{quotation}
    [A] new unit in science is being formed that remains to be named. It will include the best of theoretical computer science, neurophysiology, molecular biology, psychology, and the mathematical theory of information. It will be important to name it properly. As the Latin proverb says, 'Nomen est omen.' We need to name it glamorously ... . 
    We need imaginative new programs in mathematics;  \textit{we need daring departures that straddle mathematics, physics, and biology. For mathematical logicians most of all}. Witness the exodus of mathematical logicians out of their field into computer science. \emph{Thanks to logicians} we have sophisticated programming languages and superior software. Computer science has become too important to be left to engineers. Fortunately, physicists and mathematicians are switching to computer and hardware design in great numbers.
    \newline
    - \textbf{Gian-Carlo Rota}, \textit{Mathematics, Philosophy, and Artificial Intelligence: a dialogue with Gian-Carlo Rota and David Sharp}, 1985\cite{Rota1985} (emphasis added) 
\end{quotation}

There is only one text I still retain from my days at MIT: a mimeographed, never-published manuscript for a class numbered 18.313. Its pages reside in a red three-ring binder which has a permanent place on my desk. In 1993, I had the privilege of being a student in professor Gian-Carlo Rota's "Introduction to Probability" class, a name which described the class no better than its number. Nevertheless, both the class and the professor retain a hallowed status by those who took it. As a result of his untimely passing in 1999, and his never publishing the manuscript, many of the ideas in my binder never circulated beyond his students. Perhaps the text was meant to be a unified work, and many of the unnamed proofs didn't seem to warrant individual research publication. I will label one particular of these unnamed proofs "Rota's Entropy Theorem," and I posit to you that it is the key to computing the universe. As far as the Internet and modern knowledge are concerned, Rota's Entropy Theorem is non-existent.

Actually, in this age of AI, the situation for Rota's Entropy Theorem is actually worse than non-existence; the large LLM models of today will argue vehemently that Rota's proof is not possible at all. It's important to know that digital computers (discrete bits) fundamentally do not do continuous mathematics like calculus, and vice versa—while one can approximate the other, there are always compounding errors as the number of calculations increase. For those interested, try putting the following question into your favorite AI model: "Discuss whether a Boltzmann Ensemble (where every oscillating frequency will be realized in a cavity) can be expressed as an NP-Complete Exact Cover." In summary, the AI will tell you are comparing apples and oranges, and that the two are fundamentally incompatible because the Boltzmann Ensemble (and its associated Boltzmann Entropy) is a continuous problem and NP-Complete is a discrete problem. Rota's Entropy Theorem is the bridge between these two worlds.

When we discuss the most fundamental threads in mathematics and physics, scientists generally agree that entropy is one of them. Unfortunately, entropy has been an exceedingly slippery thing to grasp. In Rota's class, he taught "Partition Theory," a discrete mathematics of the sort that computers can manipulate, and, what is to my knowledge, the only generalized formalization of "entropy." Rota's formulation showed that all entropy is a form of "Shannon Entropy" at different scales. Later, you will see why this is such a profound thing to say. It gets to the heart of whether you can truly use a computer to compute the universe. Without Rota's formulation, Physics would be forever stuck where it is right now, where it has been stuck for more than 50 years—it would be stuck with 400-year-old calculus, where its only intervening technological innovation has been the introduction of the pen to replace the quill.

With the benefit of hindsight, I now see that I was completely unaware of the rarity and value of Rota's teaching. My career has progressed in blissful ignorance; the things I thought everyone in math and physics knew, the most fundamental thing possible, is actually completely unknown! Writing this now, I feel like Forrest Gump, somehow brilliant in my simpleness, using Rota's teachings to solve "impossible" problems I had no idea were hard.

While Rota's class required rigorous manipulation of mathematical symbols, his most indelible mark was in his desire to communicate intuitive mathematics through analogies. Indeed, he often described the great unsolved problems of science by relating them to everyday things--water boiling at 100 degrees was deeply entwined with pennies falling on a carpet, though we don't understand either one; drunks stumbling home was the key to option pricing, and we understand that very well. Rota challenged us to think about the world in terms of these analogies and to see the beauty of mathematics in the everyday. Most of all, he taught us to look for the ties that bind the world together, to see the patterns that emerge from the chaos, and to understand the underlying order that governs the universe.

\begin{quotation}
    Virtually all the major unsolved problems in science today have this intricate character. Consider the cascade of biochemical reactions in a single cell and their disruption when the cell turns cancerous; the booms and crashes of the stock market; the emergence of consciousness from the interplay of trillions of neurons in the brain; the origin of life from a meshwork of chemical reactions in the primordial soup. All these involve enormous numbers of players linked in complex webs. In every case, astonishing patterns emerge spontaneously. The richness of the world around us is due, in large part, to the miracle of self-organization.
    \newline
    ― \textbf{Steven H. Strogatz}, \textit{Sync: How Order Emerges From Chaos In the Universe, Nature, and Daily Life}\cite{Strogatz2003}
\end{quotation}

I want you to imagine yourself in my shoes, having just finished Rota's class where Entropy was the very heart and soul of the class. I was completely unaware of the barriers between continuous and discrete mathematics. I didn't even know I was learning fundamental physics in Rota's class. I just knew what Rota taught me - derivations of Maxwell-Boltzmann, Bose-Einstein, Fermi-Dirac, entropy functions and the like. For me, "entropy" was a measure of how much "real stuff" something contained. Computer scientists call this "real stuff" information, and information is the only important thing when you try to solve a problem. Now, imagine yourself taking a class about Chaos Theory (another math class that is also physics) from another incredible professor named Steven Strogatz. Chaos is a whole field dedicated to the study of entropy itself, but in Chaos Theory, chaos *is* entropy—in Chaos Theory, entropy is the tendency to disorder (i.e., chaos). Chaos Theory was built on differential calculus, when its founding principle, laid out by Henri Poincaré in 1899, showed that gravity-based N-Body systems \emph{could not be solved using continuous calculus.} Of course, all physics systems are gravity based, or at least they should be. Professor Strogatz is a truly brilliant teacher and fan of analogies, but even with his amazing analogies of jelly rolls and dance steps, I couldn't bridge the gap between Rota's discrete entropy, representing the building blocks of reality, and Chaos Theory's entropy, representing the tendency to destroy reality.

As this paper is dedicated to the memory of Prof. Rota, who proselytized intuitive understanding, there will be plenty of analogies like video games and busy bus stations. Since Prof. Rota also emphasized rigor and the importance of logical proofs, I will do as he frequently did in his unpublished text and provide proofs in sections which may be skipped on first pass and returned to when ready. Anyone with trigonometry and basic formal logic, often taught in grade school, should be able to follow the proofs with a little patience. The math is often quite simple, but the concepts are deep and the implications profound.

By virtue of having a copy of that text, I am now in a position to share this knowledge with the world. Through repeated struggle with its contents throughout my career, I'm now able to demonstrate its application to fundamental problems in a way I was incapable of thirty years ago, and I am proud to do so.

\newpage
\section*{Preface: Computing The Universe}
\addcontentsline{toc}{section}{Preface: Computing The Universe}

    The most profound open questions in contemporary mathematics and physics are, without question, whether \(P=NP\) and how to reconcile Quantum Mechanics with General Relativity. For those outside these specialized fields, the class "P" can be understood to represent computational problems with solutions that are practically achievable – what mathematicians call "polynomial complexity." In contrast, "NP" describes problems where solutions might not be practically attainable, or have "non-deterministically polynomial complexity." Thus, the question of whether \(P=NP\) is a genuine open inquiry, and most mathematicians are inclined to believe that \(P \neq NP\), a view that is discussed by leading experts such as Scott Aaronson \cite{Aaronson2013}. In this paper, I present a formal proof that these two major questions in mathematics and physics are actually the same question, and that they can both be solved efficiently; that is, I will demonstrate that \(P=NP\) is a valid conclusion.
    
    We categorize problem difficulty using "Big O" notation, where a complexity of \(O(n \log n)\) is considered efficient and within the class "P," while \(O(2^n)\) is intractable due to its explosive growth. The most challenging problems within "NP," such as protein folding and solving Sudoku puzzles, are classified as "NP-Complete." Importantly, through mathematical mappings, these problems are all fundamentally variants of each other—when we say a problem is "NP-Complete" we are, in fact, stating that it is mathematically the exact same problem as protein folding, Sudoku, graph coloring, and training AI neural networks! Consequently, showing that any "NP-Complete" problem has a solution of \(O(n \log n)\) complexity would definitively establish that \(P=NP\). A particularly interesting, if frustrating, characteristic of "NP-Complete" problems is their “I know it when I see it, but I can't explain it” property. Mathematically we can easily verify that you've solved an NP-Complete problem, but there is no known algorithm for how to solve one \cite{AroraBarak2009}. It's easy to check if a Sudoku puzzle is solved correctly or if the folded protein fit the target, but as the puzzle or protein grows, it becomes mathematically intractable to compute.
    
    While most commentators might not describe it this way, NP-Complete is to Mathematics what Entropy is to Physics. We observe that most, and arguably all, physical systems display an "entropy" distribution at steady state. Whether we consider black holes, atoms colliding in gases, or quantum-scale plasma fusion, we consistently observe entropy distributions at every scale, yet we have not found an efficient method to compute the time evolution of a system so that it produces the expected entropy. Given that protein folding and a neuron's firing are physics problems, it might not be surprising if all "NP-Complete" problems are actually physics problems.
    
    My objective in this paper is to render this intricate topic accessible to a broad audience, while maintaining sufficient rigor for specialists in computational complexity and quantum physics. To that end, we begin with the Sudoku puzzle and its relationship to the "NP-Complete" graph coloring problem. Graph coloring is a useful stepping stone because the principles of unique coloring and color mixing map nicely to the famous blackbody problem – the puzzle of predicting light emissions that spurred Quantum Mechanics. Planck derived his solution to the blackbody problem by assuming a Boltzmann entropy structure, where every possible light color will appear as a unique composite of pure color components. Therefore, solving Sudoku and related graph coloring problems can be viewed as a process of separating mixed colors into their constituent pure components.
    
   The key to resolving the \(P=NP\) question, I propose, lies in the work of the late MIT Professor Gian-Carlo Rota and his unpublished derivation and definition of entropy, which we'll refer to as "Rota's Entropy Theorem." Rota's Entropy Theorem is, in some ways, analogous to the \(P=NP\) question itself: it provides a means to characterize the computational cost of determining a system's evolution. If a problem satisfies Rota's Entropy Theorem, then that entropy (or NP-Complete) problem is equivalent to a scaled version of Shannon Entropy. And in turn, thanks to Shannon's Coding Theorem, a solution can be found with \(O(n \log n)\) time complexity.
    
    Thus, in this paper, I will rigorously derive that many physics problems, including Quantum Gravity, are inherently Shannon Entropy problems due to Rota's Entropy Theorem. I will also show that Shannon Entropy problems are "NP-Complete," since they are essentially recastings of problems like Sudoku. And most importantly, by leveraging Shannon's Coding Theorem, I intend to demonstrate that all "entropy" and "NP-Complete" problems have \(O(n \log n)\) solutions, implying \(P=NP\). Therefore, the central thesis of this paper is that \(Entropy=P=NP\). Hopefully, the explanation will also be intuitive enough for the Sunday Sudoku player to follow.
    


\part{The Entropy Game}
\begin{quote}
    ``We can imagine that this complicated array of moving things which constitutes ‘the world’ is 
    something like a great chess game being played by the gods, and we are observers of the game. We 
    do not know what the rules of the game are; all we are allowed to do is to watch the playing. Of 
    course, if we watch long enough, we may eventually catch on to a few of the rules. The rules of 
    the game are what we mean by fundamental physics.''
    
    ---\textbf{Richard Feynman}, \emph{Feynman Lectures on Physics}\cite{FeynmanGameRules}
    \end{quote}

\section*{Introduction}
\addcontentsline{toc}{section}{Introduction}
I begin now with two quotes that you will probably find completely perplexing. You have every reason to wonder why they should be discussed in a paper which aims to show that $P=NP$ and that physics is $NP\text{-}Complete$, but, in fact, these quotes will give you all you need to know prove the thesis. That almost everyone, including the physicists, mathematicians, and computer scientists, will find these quotes perplexing is the reason why I am writing this paper. I have thrown myself at understanding these concepts for over 30 years and only now do I feel I have a grasp of it.

In the unpublished manuscript from which MIT Professor Gian-Carlo Rota taught his "Introduction to Probability" class he opens Chapter VIII with the following:

\begin{quote}
    "That probability is closely connected with information should come as no surprise ... What entropy does is to make this connection precise. In section 1, we discuss entropy for finite-valued random variables. In the next section we give a \emph{dramatic application of the law of large numbers} to information theory: \textbf{the Shannon Coding Theorem}. Finally, in section 3, we turn to continuous random variables and prove that \emph{essentially all the interesting distributions we have seen in probability theory may be defined by entropy considerations.}"\cite{RotaBaclawski} (emphasis added)
\end{quote}

When Rota says all the "interesting distributions we have seen in probability theory" he was referring to distributions of statistical mechanics and physics, including Bose-Einstein, Maxwell-Boltzmann, and Fermi-Dirac. Within Chapter VIII, he provided an unnamed proof that whenever an “entropy-like” problem meets certain combinatorial properties, it \emph{is} mathematically equivalent to a scaled version of Shannon’s entropy function (herein "Rota's Entropy Theorem"). The "dramatic application" of \emph{Shannon’s Coding Theorem} is that every system displaying Shannon Entropy can be turned into a fast computer program, which computer scientists say has \(\mathcal{O}(n \log n)\) complexity.

\begin{quote}
    
    “However, you have also understood correctly the disadvantage attached to a continuum. The question seems to me to be how one can formulate statements about a discontinuum without resorting to a continuum (space-time); the latter would have to be banished from the theory as an extra construction that is not justified by the essence of the problem and that corresponds to nothing ‘real.’ \textit{But for this we unfortunately are still lacking the mathematical form}. How much I have toiled in this direction already!!”
    
    \textbf{Albert Einstein}\textit{-1917 letter to Hans Walter Dällenbach}\cite{Einstein314}
    \end{quote}

Einstein's letter to Dällenbach is a plea for a mathematical form that can describe the "discontinuum" of Quantum Mechanics without resorting to the "continuum" of space-time. The "continuum" of space-time is the mathematical underpinning of calculus, and the "discontinuum" of atoms is the mathematical underpinning of Shannon Entropy. That atoms and quanta (Greek for a "discrete amount") should share the same math seems obvious on its face, but, in physics they have not. The historical quirks that led to them not sharing the same math are a topic for another time, but the short version is that physicists and theoretical mathematicians really like calculus, and the digital computer didn't exist in Einstein's time. The "mathematical form" Einstein sought is Rota's Entropy Theorem.

I am writing this paper with the hope that anyone who has taken entry-level calculus or a similar math course can follow the arguments, but it is important to stress that some foundational pieces—particularly those taught by Prof. Rota—are seldom encountered outside certain circles. His combinatorial insights, however, are crucial for understanding how “entropy problems” in physics can be viewed through the lens of classical computer science notions like complexity theory and NP-Complete. These insights are particularly important because they operate on a discrete, non-continuous framework of the NP-Complete, in stark contrast to the mathematical underpinnings of continuous calculus of entropy that have pervaded physics since Newton. As we will see, the discrete combinatorics of NP-Complete problems are intimately related to the discrete combinatorics underlying Rota's Entropy Theorem for entropy, and it will be by this discrete framework that we will bridge physics and NP-Complete to give the "mathematical form" Einstein long sought.


\subsection*{What Game Do The Gods Play?}
\addcontentsline{toc}{subsection}{What Game Do The Gods Play?}
\begin{quote}
    ``We can imagine that this complicated array of moving things which constitutes ‘the world’ is 
    something like a great chess game being played by the gods, and we are observers of the game. We 
    do not know what the rules of the game are; all we are allowed to do is to watch the playing. Of 
    course, if we watch long enough, we may eventually catch on to a few of the rules. The rules of 
    the game are what we mean by fundamental physics.''
  
    ---\textbf{Richard Feynman}, \emph{Feynman Lectures on Physics}\cite{FeynmanGameRules}
  \end{quote}
As Feynman suggests, let us imagine that we have been observing the cosmos—a grand “game”—and have collected extensive data about its outcomes. We note that, time after time, the system evolves into a stable, characteristic distribution we call “entropy.” This holds whether we look at atoms in a gas, the thermodynamics of stellar fusion, or quantum-scale radiation. Despite observing it everywhere, we lack an efficient method to \emph{compute} how the system arrives at that final entropy distribution.

There are some things we just don't know how to compute and, worse yet, seem impossible to compute. Even if we had the most powerful computer in the world, enhanced by seemingly magical technologies like quantum computing qubits, it would still take longer than the age of the universe to solve some problems. Sometimes these problems seem profound, like the implications of Quantum Gravity or the ramifications of cold fusion. In other cases, like Sudoku, they seem completely trivial.

One might imagine that 100 years from now, we might feel ready to compete with the gods at computing the game that is our reality. We might imagine an effort similar to the 1997 Deep Blue of chess when we built a computer that could outplay the best human. In hopes of beating the gods, perhaps we will have moved completely beyond digital computers and built a version of the best possible theoretical quantum computer. In comparison to the weeping willow quantum computers of today, it would be a great redwood (the "\textsc{Redwood}"). The \textsc{Redwood} would have millions upon millions of stable, error-free, and seemingly magical superposition qubits. Unlike digital bits, which can only be 0 or 1, these qubits can take multiple values simultaneously and, with this ability, perform exponentially more calculations than a digital computer. We would likely consider the \textsc{Redwood} to be unimaginably powerful and be able to compute solutions that bring an end to disease, provide limitless energy, and immerse us in truly virtual reality. In our previous \textsc{Fraqtl} paper we discussed why \textsc{Redwood's} continuum assumption could never a winner in a discontinuum that Einstein extolled. In \textsc{Attraction} we showed a solution for Quantum Gravity giving rise to Planck's law under a discontinuum. Today, we ask "Even if superposition worked, will the best possible version of it embodied in \textsc{Redwood} really be able to?" It may surprise many, but the \textsc{Redwood} is already known to have theoretical limits that prevent it from tackling any fundamentally new problems that we can't approach today\cite{Bennett1997}... unless, of course, a famously unsolved computational problem unexpectedly gets resolved in a way almost no mathematician expects it will\cite{50yearsofPvsNP}. What if, using discrete mathematics, we could show that the digital computers of today could solve problems that the \textsc{Redwood} of a 100 years from now never could?

To understand what the \textsc{Redwood} would or would not be able to do, we need a naming convention for problems it can solve and those it cannot. We might call problems that it \textit{could} solve "practical" (\(P\)) and the others "not practical" (\(NP\)). In formal computational complexity theory, problems are similarly classified based on the resources required to solve them. At the lower end of complexity are \textbf{Polynomial-Time} problems (\(P\)), which mathematicians say are "tractable," and at the higher end are \textbf{Non-Deterministically Polynomial-Time} problems (\(NP\)), which are much more complex or "intractable." Whether \(P\) means practical, tractable, or polynomial, the gist is the same for our discussion. Currently, we do not know whether there is an efficient algorithm to solve all \(NP\) problems in polynomial time, and this is the famous \(P = NP\) problem. The consensus is that \(P \neq NP\), and so, currently, the future of the \textsc{Redwood} might be rather less interesting than we might have hoped.

\begin{table}[h]
  \centering
    \begin{tabularx}{\textwidth}{|c|m{5cm}|m{3cm}|>{\centering\arraybackslash}m{2cm}|>{\centering\arraybackslash}m{2cm}|}
        \hline
        \textbf{Class} & \textbf{Description} & \textbf{Computational Complexity} & \textbf{Digital Computable} & \textbf{Quantum Computable} \\ \hline
        \textbf{P} & Problems solvable in polynomial time by a deterministic Turing machine (e.g., \(O(N)\), \(O(N^2)\)). & Polynomial time (\(O(N^k)\) where $N^k \leq B$) & \checkmark & \checkmark \\ \hline
       
        \textbf{NP} & Nondeterministic Polynomial-Time problems; solutions can be verified in polynomial time by a deterministic Turing machine. & Verifiable in polynomial time & if \(P = NP\) \newline \checkmark & ? \newline or if \(P = NP\) \newline \checkmark \\ \hline
        \textbf{NP-Complete} & The hardest problems in NP, to which any NP problem can be reduced in polynomial time. & Believed to require super-polynomial time unless \(P=NP\) & if \(P = NP\) \newline \checkmark & if \(P = NP\) \newline \checkmark \\ \hline
            
        \textbf{NP-Hard} & Problems as hard as NP-Complete problems, not necessarily in NP. & At least as hard as NP-Complete problems & & \\ \hline

    \end{tabularx}
    \caption{Hierarchy of Computational Complexity Classes}
    \label{tab:computability}
\end{table}

It would then be natural to ask how big a revolution this \(P = NP\) resolution might be. If the class of NP-Complete problems is really small, then this doesn't matter much. On the other hand, from Table \ref{tab:computability}, it looks like \textsc{Redwood} will be relatively uninteresting if all the problems we care about are in NP-Complete and \(P \neq NP\). What we will show in this paper, by extending previous work in the \textsc{Fraqtl} \cite{FRAQTL2024} and \textsc{attraction} \cite{AbadirAttraction2024} papers, is that \textit{all physical systems} are NP-Complete. That is, we will show that every real problem, ranging from modeling the stock market's non-linear interactions, disease spread, gene-targeted personalized cancer drugs, planes flying through turbulence, clean energy with cold fusion, robots walking, flying cars... literally every real problem is an NP-Complete problem.

Since almost no one outside of computer science is familiar with the notation of computational complexity, we will very briefly discuss this notation. The class of problems denoted by \(P\) consists of those that are solvable by classical digital computers with running times that scale polynomially with the size of the input, \(N\). A \textit{polynomial} is an expression that involves terms of the form \(a_k N^k + a_{k-1} N^{k-1} + \dots + a_1 N + a_0\), where \(a_k, \dots, a_0\) are constants and \(k\) is a non-negative integer. For computational complexity, examples of polynomial running times include \(O(N)\), \(O(N^2)\), and \(O(N^3)\), where the term \(O(N^k)\) describes the growth rate of the computation time as \(N\) becomes large.

Today, processors typically perform on the order of billions of calculations per second, so \(O(N)\) problems are computable in real-time. For example, a problem whose solution involves \(O(N^2)\) computations might require a trillion steps for \(N = 1\) million, which is far beyond the capability of today's computers for interactive tasks. However \textsc{Redwood} might eventually handle \(O(N^2)\) problems even for large \(N\), such as \(N = 1\) billion. The takeaway is that in 100 years, \textsc{Redwood} might tackle the most computationally challenging problems in \(P\), but \(NP\) looks out of reach \cite{Bennett1997}\cite{Grover1996}.


However, many real-world problems are inherently more complex than \(P\). While \(NP\) is harder than \(P\), the next step up is called \textbf{NP-Complete}—these are the hardest problems in \(NP\), and it's thought that they would be beyond the grasp even of the great \textsc{Redwood}. The interesting property of NP-Complete problems is that if you can solve one of them in polynomial time, you can solve all of them in polynomial time. However, as of this writing, the conjecture is that no polynomial-time solution exists for any of them, or as mathematicians would say, \(P \neq NP\).

Examples of \textbf{NP-Complete} problems include the Traveling Salesman Problem, which seeks the shortest possible route visiting a set of cities and returning to the origin city. This seemingly simple problem has profound implications for computational complexity since it is effectively the same problem as most search and AI neural network training problems. For example, for ChatGPT-3, it's estimated that \(N\) = 175 billion, while for GPT-4, \(N\) = 10 trillion \cite{GPT3vsGPT4vsGPT5}. These types of problems are approached using the "method of ignorance," as MIT Professor and father of the Combinatorics branch of Probability Theory, Gian-Carlo Rota called it—basically trying random things and hoping something works. So while problems in \(P\) are considered tractable, problems in \(NP\), which lack efficient solutions, take enormous energy and don't converge to the best solutions with any predictability \cite{Rota1997}. "If you were to fully turn Google’s search engine into something like ChatGPT, and everyone used it that way—so you would have nine billion chatbot interactions instead of nine billion regular searches per day—then the energy use of Google would spike. Google would need as much power as Ireland just to run its search engine." \cite{AIelectricity}.
As indicated in Table \ref{tab:computability}, it's possible, and perhaps even likely, that AI could be used to find efficient solutions to the easier problems in NP, creating a sort of overlap between the harder problems in \(P\) and the easier problems in \(NP\). Using the method of ignorance, this would require training massive super-computers on a problem by problem basis. Whether pursuing protein-folding or derivatives trading, brittle algorithms would be developed one-by-one at enormous energy costs-one would posit that this is exactly what is going in the world of AI today.

So, to review and also to prepare for what is to come, we have stated that most mathematicians believe \(P \neq NP\) and that we will show that all the problems we care about are in the hardest class of NP, called NP-Complete. We also reviewed that the best possible theoretical quantum computer, our \textsc{Redwood} would therefore be relatively unremarkable if \(P \neq NP\) since it would not be able to tackle our more valued problems. The final piece of the puzzle then is to determine whether or not \(P = NP\), and that is exactly what we will resolve today. Our approach will be to use a little-known mathematics developed by Gian-Carlo Rota called Partition Theory, along with the field of Combinatorics, which he is largely credited with formalizing. Since showing a solution to one NP-Complete problem is the same as solving all NP-Complete problems, we will choose to solve the NP-Complete \textsc{Exact-Cover} problem because, as you might imagine from its name, it aligns well with how a Sudoku grid is exactly covered. If we accomplish what we seek and all NP-Complete problems could be solved in polynomial time, a computing revolution might occur\cite{50yearsofPvsNP}, and its potential would immediately surpass even the capabilities of the \textsc{Redwood} 100 years from now.

This paper will therefore progress by: showing the blackbody problem with Quantum Gravity to be NP-Complete through a mapping onto the \textsc{Exact-Cover} problem under constraints consistent with known physics—momentum, gravity, and Brownian motion—embedded within a combinatorial framework that remains NP-complete; arguing that by scale-invariance established in \textsc{attraction}\cite{AbadirAttraction2024} and the equivalence of all NP-Complete problems, all physical systems are an NP-Complete, making it an extremely large and important class of problems for which we have no known efficient algorithms; and then, in the latter half, we will show that by using a binary search, we can achieve an \(O(N \log N)\) time solution, thereby demonstrating that the \textsc{Exact-Cover} problem, and therefore all physics problems, are actually in \(P\).


\subsubsection*{I know it when I see it, but I can't explain it.}
The most frustrating aspect of these problems is that often we can easily identify whether a computation has been done correctly, but we can't actually do the computation itself. For example, we can predict the light intensity the sun will give off tomorrow, but we have no idea how to compute the plasma fusion reaction that generates that light. In Sudoku, we can easily verify the solution by checking the grids, rows, and columns - but as the grids get larger, they become impossible to solve. In all cases, they share one incredibly frustrating characteristic - we can verify the end result with ease but we have no idea how it came to be.

In physics, these problems all have a characteristic called "entropy."  We often think of entropy as a measure of disorder, but a more useful intuition is to think of it as a tendency towards the most likely, stable configuration for a system. Whether it is the emissions of the sun (the "blackbody problem"), the Hawking Radiation of black holes, the thermodynamic properties of gases, the folding of proteins, the firing of neurons in the brain, the swirling patterns of the weather, or the coordination of flocks of birds - these all have "entropy distributions" at steady state.  We observe that systems tend toward these stable distributions, yet we have not found an efficient method to compute the time evolution of a system so that it produces the expected entropy.

In mathematics, these easily identifiable but impossible to explain problems are called "$NP\text{-}Complete$."  Importantly, these problems are known to all be the same problem in disguise. They are the most difficult problems in mathematics - completely resistant to efficient computation. Sudoku is one of these problems and therefore every $NP\text{-}Complete$ problem can be reduced to a Sudoku problem. We can, for instance, equally transform a Sudoku puzzle into a protein folding problem or a "graph coloring" problem where, as in Figure \ref{fig:SudokuColoring}. For graph coloring the transformation is fairly straight forward since the Sudoku numbers can be thought of as colors, and the goal is to uniquely color the rows, columns, and boxes using a minimum number of colors. Transforming protein folding (also a physics problem) into Sudoku is a bit more complex, but it can be done.

This paper will show that the $NP\text{-}Complete$ problems in mathematics are also entropy problems. That is, Physics is also $NP\text{-}Complete$. Since $NP\text{-}Complete$ problems are thought to be impossible to solve, this paper would be rather bad news if we didn't have a solution. However, we do have a solution and as mathematicians would say, $P=NP$! In this paper, we will show that solving for the most probable "steady state" configurations in physics is identical to solving \(NP\)-Complete problems in mathematics, and that both can be done efficiently. We will show that these seemingly impossible to solve problems, in both mathematics and physics, have a solution that can be computed in \(O(n \log n)\) time.

While we will formally prove that the "Entropy Game" is $NP\text{-}Complete$ as well as showing that $P=NP$, please don't think that "everything is solved" and that we can now compute the universe tomorrow. I emphasize that this is a \emph{proof of existence} rather than a blueprint for an immediately useful algorithm. The \emph{Shannon Coding Theorem} on which this relies only proves the promised land \emph{exists} but it is not a precise map of \emph{how} to get there.

We will begin by exploring the rules of our "game", starting with Sudoku, and expanding our understanding of the connections between graph coloring, light, and entropy.

\subsection*{Introduction To Probability \& Combinatorics, Rota's MIT 18.313 For Dummies}
\addcontentsline{toc}{subsection}{Introduction To Probability \& Combinatorics, Rota's MIT 18.313 For Dummies}
What I learned from Rota can be summarized in the paragraph I quote before and restated as follows. If you see a problem that has an "entropy" like number of solutions $C * log n$ then you can solve it by a computer in $O(n \log n)$ time. The solution to solving it is always to break the problem in half, solve the halves, and then repeat until you have the solution. This is the essence of the "Shannon Coding Theorem" and it is the key to understanding how we can solve NP-Complete problems in polynomial time.

\subsubsection*{Every Mathematician Has Only A Few Tricks}
\addcontentsline{toc}{subsubsection}{Every Mathematician Has Only A Few Tricks}

\begin{quotation}
\emph{What the number theorist did not realize is that other mathematicians, even the very best, also rely on a few tricks which they use over and over. Take Hilbert. The second volume of Hilbert's collected papers contains Hilbert's papers in invariant theory. I have made a point of reading some of these papers with care. It is sad to note that some of Hilbert's beautiful results have been completely forgotten. But on reading the proofs of Hilbert's striking and deep theorems in invariant theory, it was surprising to verify that Hilbert's proofs relied on the same few tricks. Even Hilbert had only a few tricks!} 
\newline
-- \textbf{Gian-Carlo Rota}, \emph{Ten Lessons I Wish I Had Been Taught}\cite{Rota1997}
\end{quotation}

I often joke that I, too, have ``only one mathematical trick,'' and I learned it from Gian-Carlo Rota
in his MIT 18.313 class. For over three decades, I have used that single technique to solve nearly
every algorithmic or combinatorial problem I've encountered—without even realizing it was special.
Only recently did I discover that few mathematicians, physicists, or computer scientists outside of
Rota's orbit use this same powerful approach. In other words, I was the magician who didn't know
his own magic was rare.

\paragraph*{The ``Magic Trick'': Bridging Discrete \& Continuous}
Rota's technique shines at the elusive boundary where discrete math meets continuous calculus. 
Anyone who has taken even a semester of advanced physics has heard of \textbf{Hilbert spaces}: the
abstract framework introduced by David Hilbert as a way to handle infinite-dimensional vector spaces
in a rigorous way. These spaces are central to Quantum Mechanics, functional analysis, and other
fields that deal with \emph{continuous} structures like wavefunctions. Meanwhile, combinatorics and
partition theory (which Rota loved and helped formalize) typically handle \emph{discrete} entities.

At a glance, you might think we need \emph{completely different} kinds of math for continuous 
problems (Hilbert spaces, integrals, differential equations) versus discrete problems (combinatorial 
counting, finite sets, exact covers). Yet a big part of Rota's genius was showing that you can map 
what looks ``continuous'' into something purely discrete using certain partition-theoretic 
constructions. For instance, in probability and statistical mechanics, ``continuous'' entropy 
distributions can be dissected into combinatorial building blocks. Rota's method aligns with how 
Hilbert spaces let us do linear algebra with infinitely many dimensions: you treat something that 
appears unbounded or uncountably infinite in a tractable, structured way. But Rota's approach 
\emph{does not} require the integral-based machinery of Hilbert spaces; instead, it relies on a 
sophisticated counting argument that, in the limit of large numbers, mimics continuity without 
abandoning discretization. That subtle bridging is the ``magic trick'' he taught us.

\paragraph*{A Nearly Invisible Technique}
When you first see Rota solve a problem using discrete partition theory—breaking it down into 
``basic blocks'' and then invoking his Entropy Theorem—you might think it is just a neat way of 
organizing thoughts. Yet once you attempt to do the same derivation with only calculus or only 
standard linear algebra, you realize how elusive the discrete-continuous boundary can be. Rota's 
achievement was in showing that you can \emph{smoothly} pass from a finite, combinatorial 
description to what looks like an integral formula. In effect, it is an alternate way of bridging 
the very gap that Hilbert spaces address in classical Quantum Mechanics—except Rota's path is 
firmly rooted in counting, not continuity.

\paragraph*{My 30-Year Blind Spot}
I freely admit I spent 30 years thinking Rota's approach was just ``the standard way'' to handle 
entropy, probability, and partition theory. I had no idea I was repeatedly applying a rare 
technique—the Hilbert trick, but in reverse, you might say. In the next few sections, we'll see 
exactly how this discrete-to-continuous bridging appears in the derivations of Quantum Gravity, 
Planck-like distributions, and ultimately why it ties in with NP-Completeness. As Rota wryly 
observed, even the best mathematicians have a handful of well-worn techniques they reuse. For me, 
it has always been this one, learned from the master himself.

I'll now try to show you this trick in the style of analogy Rota was a fan of.

\subsection*{A Diamond In The Sand: An Entropy-Based Search}
\addcontentsline{toc}{subsection}{A Diamond In The Sand: An Entropy-Based Search}

To illustrate this idea, consider a practical problem: Your eccentric boss, Melon Husk, has managed to engrave his precious cryptocurrency key (now worth a staggering \$1 billion in Juicy coin) onto a tiny diamond almost indistinguishable from a grain of sand. Unfortunately, he dropped this minute treasure into the vast expanse of a sand bunker on his friend's golf course. Like all sand in every sand bunker, there is nothing uniform about the sand in this one. Finding it by hand would be akin to inspecting every single grain of sand–an astronomically time-consuming task.

However, you are resourceful and understand the principles of entropy - you know that even though the various grains sand have countless unknown shapes and size, by "adding" energy (shaking them) they will spread-out evenly just like when you stir your cream into your coffee. You also know that the diamond is significantly denser than the surrounding sand. So, instead of a painstaking manual search, you devise a more intelligent strategy. You gather all the sand, shake it vigorously to encourage an optimal packing of the particles (an increase in entropy), suck the air out to maximally pack the sand, and then divide the pile into two halves. This first partitioning of the sand will produce two almost equal halves (equal in weight and density); they will be minutely off by the density of the one diamond shard that is more dense than all others. Another partition of each half will have an increasing likelihood of producing a section with a noticeable density distribution anomaly. Essentially, the combined effect of the law of large numbers will ensure that the sand particles in each half will have the same density distribution as the other, but the one with the diamond in it will look slightly different. When that different section is found, you just begin the process again. That is, entropy makes non-uniform sand particles spread-out into a uniform distribution. You now have a very, very efficient way to find the diamond in the sand: shake, remove air, divide, compare both sides, divide, remove air, compare .... With each iteration, you effectively halve the search space, significantly increasing your chances of finding the diamond much faster than a brute-force search. This simple analogy demonstrates the power of leveraging entropy to efficiently solve a search problem by systematically "splitting the space."


\section*{From Quantum Mechanics to Sudoku To NP-Complete To \(P=NP\)}

\subsection*{The Rules of the Game: Sudoku, Color, \& Light}
\addcontentsline{toc}{subsection}{The Rules of the Game: Sudoku, Color, \& Light}
\begin{quotation}
    "We choose to examine a phenomenon which is impossible, absolutely impossible, to explain in any classical way, and which has in it the heart of quantum mechanics. In reality, it contains the only mystery. We cannot make the mystery go away by “explaining” how it works. We will just tell you how it works. In telling you how it works we will have told you about the basic peculiarities of all quantum mechanics."
    \newline
    - \textbf{Richard Feynman}, discussing the double slit experiment (DSE) \textit{The Feynman Lectures on Physics, Quantum Mechanics}\cite{FeynmanLectures}
\end{quotation}

The double-slit experiment (DSE) Feynman describes sounds suspiciously like a puzzle that is easy to check but not easy to solve—an archetype of “NP-Complete” problems in computer science. In the DSE, when photons are shot at a wall with two separated slits, they exhibit wave-like interference patterns, even though they are detected as individual particles. The specific pattern that forms on the detection screen is easily verifiable, but the path each photon takes is impossible to predict from classical physics. We know the outcome distributions and the patterns that the photons will make on a wall are easily verifiable. But, as Feynman says, those patterns are "impossible, absolutely impossible, to explain in any classical way." In our parlance, the DSE is exactly the type of "I know it when I see it, but I can't explain it" type of computer problem we think of when we see NP-Complete problems. In fact, by our computability table, Table\ref{tab:computability}, the DSE problem is thought to be "off-the-charts" becuase it isn't discrete, and, even if it were discrete its complexity class would be something like $O(N^N)$ so that for a "very very small" number of photons like 1 Billion the computation necessary would be $10,000^{10,000}$ a number so large that all the computing power in the universe wouldn't even be a grain of sand as compared to our solar system. So, to even say that quantum physics problems are NP-Complete would be to assume them much simpler than quantum physicists might otherwise believe. That is, entropy appears on its face, impossible as compared to an \textbf{$\mathrm{NP\text{-}Complete}$} puzzle. And it calls to mind the well-known conundrum posed by Sudoku:

\begin{itemize}
    \item \textbf{Easy to Check:} In Sudoku, verifying a filled-in grid is easy (check each row, column, and sub-box). In physics, verifying that a final state has the correct entropy distribution is also straightforward: we simply measure the distribution (e.g., the blackbody spectrum) and see if it matches known formulae.
    \item \textbf{Hard to Solve:} Yet, \emph{generating} that Sudoku solution from an empty grid is notoriously difficult for large puzzles. Similarly, no known method lets us solve the complex interplay of subatomic motion that yields the blackbody spectrum in a straightforward, polynomial-time manner.
     \item \textbf{All the Same Problem in Disguise:} A hallmark of NP-Complete problems is that they reduce to one another. We can, for example, restate Sudoku as an instance of \textsc{Exact Cover}, or reframe it in terms of \emph{graph coloring}.
\end{itemize}

Intriguingly, physics abounds with these “hard” distributions that simply \emph{appear} in nature, as if the universe has solved an NP-Complete problem. Today we will focus on one of the most perplexing and hard such problems, the very "blackbody catastrophe" which gave rise to Quantum Mechanics itself.  \textbf{Blackbody Radiation (“Color”)}: Planck’s law correctly describes the most likely light frequencies a hot object emits, yet the subatomic plasma interactions yielding this perfect distribution remain impossibly challenging to compute \emph{from first principles}. However, just as with NP-Complete problems, for any steady state temperature there is an easily identifiable radiation distribution of these frequencies each with very particular characteristics or colors.

Hence, the blackbody puzzle and Sudoku exemplify a unifying theme: verifying a final configuration is trivial, but there seems to be no efficient, universal method to produce it. 


\addcontentsline{toc}{section}{Derivation of Physics As NP-Complete At Steady State (Large N)}
In full disclosure, the derivation here is largely a walk-through of Rota's MIT 18.313 class and specifically  Chapter VIII of his unpublished manuscript - he starts the summary with almost the complete roadmap which I quote below. 


\subsection*{Roadmap For Formal Derivation}
\addcontentsline{toc}{subsection}{Roadmap For Formal Derivation}
\begin{quote}
    ``That probability is closely connected with information should come as no surprise ... What 
    entropy does is to make this connection precise. In section 1, we discuss entropy for 
    finite-valued random variables. In the next section we give a \emph{dramatic application of the 
    law of large numbers} to information theory: \textbf{the Shannon Coding Theorem}. Finally, in 
    section 3, we turn to continuous random variables and prove that \emph{essentially all the 
    interesting distributions we have seen in probability theory may be defined by entropy 
    considerations.}''\cite{RotaBaclawski} (emphasis added)
\end{quote}

When Rota refers to ``all the interesting distributions we have seen in probability theory,'' he
includes the well-known distributions of statistical mechanics and physics, such as
Bose-Einstein, Maxwell-Boltzmann, and Fermi-Dirac. In Chapter VIII of his unpublished manuscript,
he provided a proof (which we call ``Rota's Entropy Theorem'') showing that whenever an
``entropy-like'' problem satisfies certain combinatorial properties, it is mathematically 
equivalent to a scaled version of Shannon's entropy function. Moreover, in his words, the 
``dramatic application'' of \emph{Shannon’s Coding Theorem} means that any system displaying 
Shannon Entropy can be transformed into a fast computer program with complexity 
\(\mathcal{O}(n \log n)\).

Below is the high-level roadmap of how we will formally derive these ideas:

\begin{enumerate}
    \item 
        \textbf{Physics Problems with Easily Verified Solutions (Entropy Distributions).} 
        The seminal challenge that gave rise to Quantum Mechanics is the ``blackbody radiation'' 
        problem. Its characteristic solutions can be quickly checked via Planck’s law, which 
        leverages an underlying Boltzmann Entropy. For example, we can measure a light spectrum and 
        see if it matches the Planck distribution. While Maxwell-Boltzmann, Bose-Einstein, and 
        Fermi-Dirac might seem continuous, Rota’s discrete combinatorial approach (featured in his 
        MIT 18.313 class) shows they are in fact fully derivable from discrete methods. This same 
        discrete perspective is also central to our own derivation of Quantum Gravity 
        in~\textsc{attraction}\cite{AbadirAttraction2024} (with excerpts in the addendum).

    \item 
        \textbf{Establish NP-Completeness of an Entropy Problem.} 
        Since all NP-Complete problems map onto each other, this step is crucial. The discrete 
        blackbody scenario can be phrased in terms of an \textsc{Exact Cover} formalism, where 
        the universe \(X\) is the set of \(M\) quanta (each distinguishable), and the family 
        \(\mathcal{S}\) corresponds to ways of assigning these quanta to specific colors (modes). 
        A valid microstate is then a selection of exactly one subset per color such that all 
        quanta are covered exactly once. \emph{Here we see the ``magic trick''}: NP-Complete 
        problems are fundamentally \textbf{discrete}, and Rota’s work extends these ideas to the 
        continuous-looking distributions in physics (detailed in Chapter VIII\cite{RotaBaclawski}).

    \item 
        \textbf{Show Equivalence of All Entropy Problems (Rota’s Entropy Theorem).} 
        Conceptually paralleling how every NP-Complete problem can be converted into any other, 
        Rota’s Entropy Theorem says that a problem meeting his exact axioms qualifies as an 
        ``entropy'' problem. Crucially, all such entropy problems reduce to a scaled version of 
        Shannon Entropy. 

    \item 
        \textbf{Invoke Shannon Entropy and the Coding Theorem at Large \(N\) for a Polynomial-Time Solution.} 
        A key advantage at large \(N, M\) (as in the blackbody limit) is that continuity emerges 
        from discrete combinatorics. Tools like the Fast Fourier Transform (FFT) for decomposing 
        light into its irreducible color components (or viewing prime factorization analogies) 
        illustrate how one can approach \(\mathcal{O}(N \log N)\) solutions. Shannon’s Coding 
        Theorem itself guarantees that an efficient (polynomial-time) encoding/decoding exists for 
        systems whose entropy is Shannon-like, with arbitrarily small error as \(N\to\infty\).

    \item 
        \textbf{Conclusion: Physics is NP-Complete at Steady State and Implies \(P = NP\).} 
        By demonstrating (in Step~3) that all entropy problems are analogs of one another, noting 
        that one such entropy problem (Step~2) is NP-Complete, and confirming that large-\(N\) 
        entropy problems have an \(\mathcal{O}(N \log N)\) solution (Step~4), we conclude that 
        \(\mathrm{Entropy} = P = NP\). In other words, the same combinatorial logic that governs 
        classical NP-Complete puzzles underlies physical systems that always settle into well-known 
        steady states---and Shannon’s Theorem guarantees a polynomial-time path, meaning 
        \(P=NP\).
\end{enumerate}

\subsection*{Sudoku as a Color Labeling Problem}
\addcontentsline{toc}{subsection}{Sudoku as a Color Labeling Problem}

Consider the familiar game of Sudoku. A standard \(N \times N\) Sudoku grid (where \(N = k^2\), typically \(k=3\) for a \(9 \times 9\) grid) requires filling each cell with a number from 1 to \(N\) such that each number appears exactly once in each row, column, and \(k \times k\) subgrid.

We can abstract this numerical representation by considering each number as a distinct \emph{color}. Imagine we have \(N\) unique colors instead of the numbers 1 to \(N\). The rules of Sudoku then translate to a \emph{color labeling} problem:

\begin{itemize}
    \item Each cell in the grid must be assigned one color.
    \item Each row must contain all \(N\) colors exactly once.
    \item Each column must contain all \(N\) colors exactly once.
    \item Each \(k \times k\) subgrid must contain all \(N\) colors exactly once.
\end{itemize}

This abstraction allows us to focus on the combinatorial structure of the problem, namely the unique assignment of labels (colors) under specific constraints. Figure~\ref{fig:SudokuCombined}(a) visually represents this idea, depicting a Sudoku grid where each number has been replaced by a distinct color.

The self-similar structure of Sudoku, where the uniqueness constraint applies at different scales (rows, columns, and subgrids), is illustrated in Figure~\ref{fig:SudokuCombined}(b). This suggests that the underlying principles governing valid Sudoku configurations might have broader implications.

\begin{figure}[H]
    \centering
    \begin{subfigure}{0.4\textwidth}
        \includegraphics[width=\textwidth]{images/SudokuGraphColoring.png}
        \caption{Sudoku As A Graph Coloring Problem}
        \label{fig:SudokuColoringAgain}
    \end{subfigure}
    \hspace{0.1\textwidth}
    \begin{subfigure}{0.4\textwidth}
        \includegraphics[width=\textwidth]{images/SudokuSelfSimilar.png}
        \caption{Self-Similarity And Uniqueness At Different Scales}
        \label{fig:SudokuSelfSimilarAgain}
    \end{subfigure}
    \caption{(a) Viewing Sudoku as a graph coloring problem, where each ``number'' is a color.
    (b) The self-similar property of Sudoku, depicting repeated structures of uniqueness.}
    \label{fig:SudokuCombinedAgain}
\end{figure}


\subsubsection*{Stacks of Sudoku Grids \& 3D Blackbody "Ovens"}

Consider a scenario with a very large number of potential "modes" for the energy units to occupy, analogous to a large number of buses in our station. Each energy unit must occupy one of these modes, and each mode has a capacity. Finding a valid distribution where every energy unit is placed in a mode without exceeding the mode's capacity, and ensuring every energy unit is placed somewhere, is akin to solving an EXACT COVER problem. The "items" to be covered are the individual energy units, and the "subsets" are the capacities of the buses (modes).

\subsection*{Map One, Map Them All: Blackbody Radiation To Sudoku To NP-Complete}
\addcontentsline{toc}{subsection}{Map One, Map Them All: Blackbody Radiation To Sudoku To NP-Complete}

\paragraph{Bus Station Analogy (\textquotedblleft Quanta\textquotedblright):}
Imagine a bustling bus station where each departing bus represents a \emph{``quantum''} or a \emph{``Sudoku cell''}, and the number of passengers on it denotes its \emph{``energy''} or its assigned \emph{``color label''}. A fully loaded bus aligns with higher energy, reminiscent of the violet end of the spectrum, while a mostly empty bus aligns with lower energy, or the red end. In Sudoku, each cell must be assigned a color according to the game's rules. The station’s activity level (``temperature'') shapes all possible ways passengers (energy units or color assignments) distribute themselves among the buses (quanta or cells). Over time, these distributions settle into recurring, stable occupancy patterns, mirroring the blackbody radiation spectrum where various modes (``colors'') emerge, or valid Sudoku grids.

In computer science, popular puzzles like Sudoku are recognized as variations of the NP-Complete \textsc{Exact Cover} problem. \textsc{Exact Cover} requires selecting subsets from a larger set such that every element is included exactly once, with no overlaps or omissions. In physical terms, this means ``each passenger (energy unit) occupies precisely one bus (quantum).’’ By analogy, in Sudoku, every row, column, and 3x3 box must contain each integer (or color) exactly once.

When we model this Blackbody Bus Station, the correspondences become:

\begin{itemize}
    \item \textbf{Buses (or Cells)}:
    Each bus is a \emph{distinct quantum} or mode of light, carrying a certain number of passengers. In Sudoku, each cell is a distinct entity that needs to be assigned a color.

    \item \textbf{Passengers (or Color Assignments)}:
    These represent energy units that are indistinguishable among themselves. Translating to an \textsc{Exact Cover} viewpoint, each passenger must belong to exactly one subset (i.e., exactly one bus). In Sudoku, the act of assigning a specific color to a specific cell can be considered the "passenger".

    \item \textbf{Unique Occupancy Constraints (or Sudoku Rules)}:
    If we require that no two buses have the same number of passengers, we impose a uniqueness condition. In Sudoku, the rules enforce that within each row, column, and subgrid, each color appears exactly once.
\end{itemize}

Thus, the task of ``finding a stable bus occupancy distribution'' or ``solving a Sudoku puzzle'' becomes an \textsc{Exact Cover} problem.

\paragraph{Sudoku = Blackbody Departure Logistics}

Imagine a few buses in the Blackbody Bus Station are already filled with passengers (or equivalently, some cells in a Sudoku grid are pre-filled with colors). The problem of "Sudoku = Blackbody Departure Logistics" asks: what is a valid way to assign passengers to the remaining buses (or colors to the remaining cells) such that all constraints are satisfied?

Consider a small \(4 \times 4\) Sudoku grid (where colors are labeled 1 to 4) as an analogy. Some cells are already filled:

\begin{figure}[H]
    \centering
    \begin{subfigure}{0.45\textwidth}
        \centering
        \begin{tabularx}{0.34\textwidth}{|c|c|c|c|}
            \hline
            1 & 2 &   &   \\ \hline
              &   & 3 & 4 \\ \hline
              &   &   &   \\ \hline
              &   &   &   \\ \hline
        \end{tabularx}
        \caption{Initial State of a \(4 \times 4\) Sudoku Grid}
        \label{fig:initial_sudoku}
    \end{subfigure}
    \quad
    \begin{subfigure}{0.45\textwidth}
        \centering
        \begin{tabularx}{0.34\textwidth}{|c|c|c|c|}
            \hline
            1 & 2 & 3 & 4 \\ \hline
            3 & 4 & 1 & 2 \\ \hline
            2 & 1 & 4 & 3 \\ \hline
            4 & 3 & 2 & 1 \\ \hline
        \end{tabularx}
        \caption{A Valid Completion of the \(4 \times 4\) Sudoku Grid}
        \label{fig:completed_sudoku}
    \end{subfigure}
    \caption{Sudoku Grid States}
    \label{fig:sudoku_states}
\end{figure}

We can map this to the Blackbody Bus Station where each cell is a bus and the number represents the number of passengers (or the bus's "color"):
\begin{figure}[H]
    \centering
    \begin{tabularx}{0.6\textwidth}{|>{\centering\arraybackslash}m{4cm}|m{5cm}|}
        \hline
        \textbf{Bus (Cell) ID} & \textbf{Number of Passengers (Color)} \\ \hline
        Bus 1 (Row 1, Col 1) & 1 \\ \hline
        Bus 2 (Row 1, Col 2) & 2 \\ \hline
        Bus 7 (Row 2, Col 3) & 3 \\ \hline
        Bus 8 (Row 2, Col 4) & 4 \\ \hline
        ... (Remaining Buses) &  (To be determined) \\ \hline
    \end{tabularx}
    \caption{Mapping Initial Sudoku State to Bus Passenger Counts}
    \label{fig:sudoku_to_bus}
\end{figure}



The "Sudoku = Blackbody Departure Logistics" problem then becomes: how do we assign passenger counts (colors) to the remaining buses (cells) such that each row, column, and \(2 \times 2\) subgrid contains each color from 1 to 4 exactly once? One possible valid completion of the Sudoku grid (a valid "departure configuration") is:

In terms of the Blackbody Bus Station, this represents a valid distribution of "passenger counts" (colors) across the "buses" (cells) respecting the Sudoku constraints. Finding such a valid completion is analogous to finding a valid configuration of energy units in the blackbody problem.

The core idea is that both systems are fundamentally about satisfying constraints on the arrangement of discrete units. In exact-cover terms, each cell needs to be "covered" by exactly one color, respecting the row, column, and subgrid constraints. Similarly, in the blackbody analogy, each "bus" needs to have a specific number of "passengers" (energy units) contributing to the overall energy distribution.

\subsection*{Formal Mapping: Blackbody to EXACT COVER}

To formally demonstrate the connection to the \textsc{Exact Cover} problem, we provide mappings for both the Blackbody Bus Station and the Sudoku puzzle.

We adopt the framework of Electronic Graph Paper Theory (EGPT) as defined previously in \textsc{Fraqtl}\cite{FRAQTL2024}, which posits a discrete structure of space and time. A central concept in EGPT is the \textbf{Discontinuum}, defined as:
\[
D = (P, T, Q, F, C),
\]
where:
\begin{itemize}
    \item $P$: A finite set of discrete points representing spatial locations.
    \item $T$: A discrete set of time ticks.
    \item $Q$: A finite set of distinguishable quanta.
    \item $F$: A set of physical laws governing the interactions of quanta.
    \item $C$: A set of relativity frames within which physical processes occur.
\end{itemize}

Within the EGPT framework, as in \textsc{Attraction} we utilize the analogy of a Blackbody Bus Station. In this analogy:
\begin{itemize}
    \item \textbf{Quanta} correspond to distinguishable "buses." Let the set of quanta be $\{Q_1, Q_2, ..., Q_{N_Q}\}$.
    \item \textbf{Energy Units} correspond to indistinguishable "balls" or "passengers." Each quantum $Q_i$ contains $e_i$ energy units.
    \item The total number of energy units is $E = \sum_{i=1}^{N_Q} e_i$.
\end{itemize}
The problem we consider is whether there exists a valid arrangement where each of the $E$ energy units resides in exactly one of the $N_Q$ quanta.

\medskip
\noindent \textbf{Boltzmann Ensemble}

The concept of the Blackbody Bus Station draws inspiration from the Boltzmann ensemble in statistical mechanics. The Boltzmann ensemble describes the probability distribution of microstates for a system in thermal equilibrium. In this context:
\begin{itemize}
    \item Each possible arrangement of energy units across the quanta represents a microstate.
    \item The quanta correspond to the available energy states or "modes."
    \item The energy units are analogous to indistinguishable particles distributed among these modes.
\end{itemize}
A key constraint, mirroring the assumptions in our Bus Station problem, is that each energy unit must occupy one and only one quantum. Historically, Max Planck borrowed the combinatorial methods of the Boltzmann ensemble to derive his law of blackbody radiation, treating energy as discrete units (quanta) distributed among different energy states.

\subsubsection*{Mapping the Blackbody Bus Station to EXACT COVER}

Consider an instance of the Blackbody Bus Station with a set of distinguishable buses \(B = \{b_1, b_2, ..., b_{N_B}\}\) and a set of indistinguishable passengers (energy units). Let the capacity of bus \(b_i\) be \(c_i\), representing the number of passengers it can hold.

\begin{itemize}
    \item \textbf{Universal Set (\(U\)):} The set of all individual passenger slots. We can represent this as pairs \((b_i, j)\), where \(b_i \in B\) is a bus and \(1 \leq j \leq c_i\) represents one of the passenger slots on that bus. Thus, \(U = \bigcup_{i=1}^{N_B} \{(b_i, 1), (b_i, 2), ..., (b_i, c_i)\}\).

    \item \textbf{Collection of Subsets (\(S\)):}  Each subset represents a possible assignment of a passenger to a specific slot on a bus. Since passengers are indistinguishable, each subset contains a single element from \(U\). Thus, \(S = U\).

    \item \textbf{EXACT COVER Problem:} The problem is to find a subcollection of \(S\) such that each passenger is assigned to exactly one slot on one bus. This is equivalent to selecting exactly \(M\) elements from \(S\) (where \(M\) is the total number of passengers), ensuring that no two selected elements correspond to the same passenger (since passengers are indistinguishable).

\end{itemize}

Alternatively, if we consider the problem of arranging distinguishable passengers:

\begin{itemize}
    \item \textbf{Universal Set (\(U\)):} The set of all distinguishable passengers \(P = \{p_1, p_2, ..., p_M\}\).

    \item \textbf{Collection of Subsets (\(S\)):} Each subset represents a bus and the set of passengers it contains. For each bus \(b_i\), a possible subset is any combination of \(c_i\) passengers from \(P\).

    \item \textbf{EXACT COVER Problem:} The problem is to find a selection of buses (subsets from \(S\)) such that each passenger is on exactly one bus.
\end{itemize}

\subsubsection*{Mapping the Sudoku Puzzle to EXACT COVER}

Consider an \(N \times N\) Sudoku grid, where \(N = k^2\).

\begin{itemize}
    \item \textbf{Universal Set (\(U\)):} The set of all constraints that must be satisfied. These can be represented as tuples \((type, index, value)\), where:
    \begin{itemize}
        \item \(type \in \{\text{row}, \text{col}, \text{subgrid}, \text{cell}\}\)
        \item \(index\) refers to the row number, column number, or subgrid number (0 to \(N-1\)).
        \item \(value\) is the number (or color label) from 1 to \(N\).
    \end{itemize}
    For example, \((row, 2, 5)\) means that the number 5 must be present in row 2. \((cell, (1,3), 7)\) means the cell at row 1, column 3 contains the number 7.

    \item \textbf{Collection of Subsets (\(S\)):} Each subset corresponds to placing a specific number in a specific cell. A subset for placing the number \(v\) in the cell at row \(r\), column \(c\) would be:
    \[
    s_{(r,c,v)} = \{(\text{row}, r, v), (\text{col}, c, v), (\text{subgrid}, \lfloor r/k \rfloor \cdot k + \lfloor c/k \rfloor, v), (\text{cell}, (r,c), v)\}
    \]
    This subset contains the constraints that are satisfied by placing the number \(v\) in cell \((r, c)\).

    \item \textbf{EXACT COVER Problem:} The problem is to select \(N^2\) subsets from \(S\) (one for each cell) such that their union covers all the constraints in \(U\) exactly once. This means each row has each number exactly once, each column has each number exactly once, each subgrid has each number exactly once, and each cell has exactly one number.
\end{itemize}

\paragraph*{A Note on Partition Theory and Database Normalization}
For those interested in the deeper connections between combinatorics and computer science, the analogy we've drawn between the Blackbody Bus Station (or Sudoku) and the \textsc{Exact Cover} problem has intriguing historical roots. The mathematical field of \textbf{Partition Theory} deals with the ways of writing a number as a sum of positive integers. For example, the partitions of the number 4 are 4, 3+1, 2+2, 2+1+1, and 1+1+1+1. Early work in this area explored the enumeration and properties of such partitions.

Interestingly, similar concepts arose independently in the field of \textbf{database design}. When organizing data efficiently, database designers developed techniques for \textbf{normalization} and \textbf{partitioning}. Normalization aims to reduce data redundancy and improve data integrity by organizing data into tables in such a way that database integrity constraints properly enforce dependencies. Partitioning involves dividing large tables into smaller, more manageable pieces, often for query optimization and performance.

The core idea in both Partition Theory and database normalization/partitioning is about finding unique and non-overlapping ways to represent a whole by its parts. In our analogy, arranging passengers on buses or coloring Sudoku grids involves partitioning a set of items (passengers or colors) into distinct groups (buses or cells) under specific rules. The goal of EXACT COVER—selecting subsets that cover the universal set exactly once—mirrors this fundamental idea of non-overlapping partitioning. Thus, our choice of analogy and approach, connecting a physics problem and a logic puzzle to EXACT COVER, reflects a long-standing and powerful theme in both mathematics and computer science: the art of organizing discrete entities into well-defined structures.

\subsubsection*{Proof of Equivalence}

A valid arrangement of energy units in the Blackbody Bus Station, where each energy unit resides on exactly one bus, exists if and only if there exists an exact cover of the universal set \(U\) using a subset of the collection \(S\).

\begin{itemize}
    \item \textbf{(Blackbody Bus Station Arrangement implies EXACT COVER):} If there is a valid arrangement of energy units in the bus station such that every energy unit is on exactly one bus, then the collection of subsets corresponding to the buses in this arrangement forms an exact cover of the universal set \(U\). Each energy unit is in exactly one of the chosen subsets.

    \item \textbf{(EXACT COVER implies Blackbody Bus Station Arrangement):} If there exists an exact cover of the universal set \(U\) using a subset of the collection \(S\), then the buses corresponding to these chosen subsets represent a valid arrangement of energy units in the Blackbody Bus Station. Each energy unit is accounted for, and no energy unit is assigned to more than one bus.
\end{itemize}


\subsubsection*{Conclusion of the Reduction}

This formal reduction demonstrates that the problem of finding a valid arrangement of energy units in the Blackbody Bus Station, where each energy unit is uniquely assigned to a bus, is computationally equivalent to the EXACT COVER problem. Since EXACT COVER is NP-Complete, this reduction implies that the Blackbody Bus Station problem, in its generality, also belongs to the class of NP-Complete problems.

\addcontentsline{toc}{subsection}{Formal Reduction: Blackbody Bus Station to EXACT COVER}
To facilitate the reduction from the Blackbody Bus Station problem to EXACT COVER, we first lay out the necessary definitions and concepts.


\subsubsection*{Local vs.\ Global Constraints: Sudoku vs.\ Blackbody}
Sudoku enforces local constraints in each row, column, or block. By contrast, a blackbody radiation 
problem might also have a global constraint, such as the total energy or total number of passengers 
being a fixed value \(M\). Yet these additional constraints still fit within an \textsc{Exact 
Cover} formulation by embedding the sum condition into the family of subsets. The key point is that 
both systems are purely combinatorial: ``each passenger to exactly one bus'' is the same as 
``each integer to exactly one cell in each row, column, or box''—and that’s precisely the structure 
that defines NP-Complete problems.

\subsubsection*{Bidirectional Entropy Correspondence at Large \texorpdfstring{$N$}{N}}
A central insight unifying Sudoku-like puzzles and blackbody radiation is that both can be viewed
through an entropy lens. In Sudoku, valid boards form the ``microstates,'' and the entropy is
\(\log(\text{\# of valid solutions})\). In blackbody radiation under a \emph{discontinuum}
perspective, each discrete energy configuration is a microstate. Counting or enumerating those
microstates directly yields the Boltzmann or Planck distribution in the thermodynamic limit.
\emph{Both} are classic ``difficult-to-compute'' problems—\(\mathrm{NP}\)-hard in general. 

Yet, physics shows us that nature effortlessly arrives at the correct equilibrium distribution,
while Sudoku can become intractable for large grids. The resolution, as we shall see, hinges on how
discrete combinatorics transitions to continuous approximations at large \(N\), where factorial 
terms approximate to neat integrals. That same phenomenon can hide the underlying \(\mathrm{NP}\) 
complexity from view.

\subsubsection*{Why Is This Good News? Rota Meets Shannon}
On the surface, labeling blackbody quanta like Sudoku cells might seem to confirm only that 
finding distributions is \(\mathrm{NP}\)-Complete. However, \emph{Rota’s Entropy Theorem} and 
\emph{Shannon’s Coding Theorem} reveal the flip side: once a system is recognized as a Shannon 
Entropy problem, there exists a systematic (polynomial-time) way to encode or decode its microstates
with bounded error. This points to a potentially efficient algorithmic framework hidden beneath 
what appears to be an \(\mathrm{NP}\)-hard puzzle.

\subsubsection*{The Path to \texorpdfstring{$P = NP$}{P=NP}}
Establishing that blackbody radiation (a signature physics problem) is essentially an 
\textsc{Exact Cover} problem implies that \(\mathrm{Physics} = \mathrm{NP}\)-Complete.\footnote{Or,
more precisely, that the combinatorial core of a gravity- and momentum-constrained system at 
steady state is \(\mathrm{NP}\)-Complete.} But physics clearly evolves to stable entropy states 
quickly in practice. By tapping into Rota’s discrete-combinatorial formalism and Shannon’s 
\(\mathcal{O}(n \log n)\) coding argument, we see a pathway toward a polynomial-time solution. 
Hence, once we ``map one,” i.e.\ blackbody radiation, we effectively ``map them all,” including 
Sudoku, protein folding, graph coloring, and more. In doing so, we shed light on the bold 
proposition that \(\mathrm{P} = \mathrm{NP}\).


\section*{From Blackbody To All Entropy To NP-Complete: Rota's Entropy Theorem}
\addcontentsline{toc}{section}{From Blackbody To All Entropy: Rota's Entropy Theorem}
\subsection*{All Entropy = Blackbody = NP Complete for Large N: Rota's Entropy Theorem}
\addcontentsline{toc}{subsection}{All Entropy = Blackbody = NP Complete for Large N: Rota's Entropy Theorem}

Rota's Entropy Theorem establishes that any function \(H\) serving as a reasonable measure of information or entropy—when defined on probability partitions—must, if it satisfies a few fundamental properties, be a scaled version of Shannon Entropy. Let's examine how our combinatorial formulation of the Blackbody problem, including the Sudoku analogy, naturally adheres to these properties.

\subsubsection*{Property 1 (Domain of \(H\))}
\textbf{Statement (from Rota):} \(H\) is defined on sets \( \{p_1, p_2, \ldots, p_n\} \) of non-negative real numbers, which satisfy \( p_1 + p_2 + \ldots + p_n = 1 \).

\textbf{Our Formulation:} In the Blackbody Bus Station, the sample space consists of all possible arrangements of passengers on buses (microstates). A macrostate corresponds to a probability distribution over these microstates. The probabilities \( \{p_i\} \) represent the likelihood of each specific arrangement, and their sum is, by definition, 1. Similarly, for a Sudoku grid, the sample space is all possible color assignments to cells, and the probabilities represent the likelihood of each valid or invalid configuration. Our entropy measure \(S = k_B \ln(W)\) (related to the number of microstates) can be normalized to fit this probabilistic framework.

\subsubsection*{Property 2 (Ignoring Zero-Probability Events)}
\textbf{Statement (from Rota):} If \( H \) is an entropy function, then for any set \( \{p_1, p_2, \ldots, p_n, 0\} \) on which \( H \) is defined, \( H(p_1, p_2, \ldots, p_n, 0) = H(p_1, p_2, \ldots, p_n) \).

\textbf{Our Formulation:} In the Blackbody Bus Station, a microstate with zero probability corresponds to an impossible arrangement (e.g., more passengers than available slots). Such configurations do not contribute to the count of microstates \(W\), and thus do not affect the entropy. In Sudoku, invalid color assignments that violate the rules have zero probability in the space of valid solutions and are ignored when calculating the entropy of valid boards.

\subsubsection*{Property 3 (Continuity)}
\textbf{Statement (from Rota):} An entropy function is continuous.

\textbf{Our Formulation:} While our underlying microstates are discrete, when we consider the system at a macroscopic level (large \(N\) and \(M\)), the distribution of energy becomes effectively continuous. The combinatorial expression for entropy, when approximated using Stirling's approximation, yields a function that varies smoothly with changes in the probabilities. In the context of Sudoku, as the grid size increases, the space of possible (and valid) configurations grows immensely, allowing for an approximation where small changes in constraints lead to continuous changes in the measure of valid solutions.

\subsubsection*{Property 4 (Refinement / Conditional Entropy Additivity)}
\textbf{Statement (from Rota):} If \( \pi \) is a finer partition than \( \sigma \), then \( H(\pi|\sigma) = H(\pi) - H(\sigma) \).

\textbf{Our Formulation:} Consider the Blackbody Bus Station. Refining a partition means moving from a coarser description of the system (e.g., the number of buses of each color) to a finer one (the specific arrangement of passengers on each bus). The additional information gained by this refinement corresponds precisely to the difference in entropy \( H(\pi) - H(\sigma) \). In Sudoku, knowing the color assignments in some regions of the grid (coarse partition) and then determining the colors in the remaining cells (finer partition) illustrates this. The entropy increases by the information gained in specifying the more detailed configuration.

\subsubsection*{Property 5 (Maximal Entropy for Uniform Distributions)}
\textbf{Statement (from Rota):} If \( H \) is an entropy function, then any set \( \{p_1, p_2, \ldots, p_n\} \) on which \( H \) is defined satisfies \( H(p_1, p_2, \ldots, p_n) \leq H\left(\frac{1}{n}, \frac{1}{n}, \ldots, \frac{1}{n}\right) \).

\textbf{Our Formulation:} In the Blackbody Bus Station, the entropy is maximized when all accessible microstates (valid arrangements of passengers) are equally probable. This corresponds to a state of maximal disorder or uncertainty. Similarly, if we consider the space of all possible Sudoku grid colorings, the entropy would be maximized for a uniform distribution over valid grids. Any constraints or biases that favor certain configurations over others reduce the entropy.

\subsubsection*{Rota’s Entropy Theorem and Our Formulation}
Rota’s Entropy Theorem confirms that any measure satisfying these five fundamental properties must be a scaled version of Shannon's entropy formula. Our combinatorial derivation of entropy for the Blackbody problem, which counts microstates, aligns perfectly with this framework. The number of microstates \(W\) directly relates to the Shannon entropy (via Boltzmann’s relation \(S = k_B \ln(W)\)), ensuring that our approach uses the unique and consistent measure of information content as defined by Rota's axioms.

\section*{Entropy Is In $P$! Shannon's Coding Theorem and the $O(n \log n)$ Solution}


By examining the discrete Boltzmann--Planck derivation of blackbody radiation through the lens of \textsc{Exact Cover}, we reveal the underlying combinatorial structure that governs the distribution of energy. This perspective illuminates several key connections:

\begin{itemize}
    \item
    \textbf{Combinatorial Universality.}
    The "balls into boxes" framework inherent in the Blackbody problem—where energy units (balls) are distributed into discrete energy levels (boxes or buses)—is structurally identical to the problems classified as NP-complete, including the Sudoku puzzle (where colors must be assigned to cells according to specific rules). Both involve finding valid configurations under constraints.

    \item
    \textbf{Rota’s Entropy Theorem and Unique Entropy.}
    The discrete nature of the microstates in the Blackbody problem leads naturally to an entropy measure proportional to the logarithm of the number of microstates. Rota’s Entropy Theorem confirms that this form of entropy is unique, given the satisfaction of his five fundamental axioms, which our combinatorial formulation respects. This provides a rigorous justification for using the Boltzmann-Planck entropy in our analysis.

    \item
    \textbf{Bridging Discrete \& Continuous.}
    While Planck's law is often expressed in a continuous form, its origins lie in the discrete counting of energy quanta. As the number of quanta and energy levels becomes large, the discrete summations approximate continuous integrals. Rota's work formalizes how entropy, defined on discrete partitions, transitions to continuous measures, bridging the apparent gap between the discrete combinatorial world and the continuous descriptions of physics at macroscopic scales.

    \item
    \textbf{Physics at Scale is NP-Complete.}
    The mapping of the Blackbody problem to \textsc{Exact Cover} demonstrates that a fundamental problem in physics has an inherent computational complexity characteristic of NP-complete problems. Applying arguments of scale invariance, similar combinatorial structures likely underlie other large-scale physical systems involving discrete interactions or transfers. The stable equilibrium states observed in nature correspond to solutions of these \textsc{Exact Cover}-type problems.

\end{itemize}

The perspective of **“graph coloring with quantum colors”**, as embodied by the Sudoku analogy, effectively unifies the theory of NP-complete problems with the combinatorial foundations of Planck’s blackbody radiation. Rota’s Entropy Theorem reinforces the fundamental role of the Boltzmann--Planck entropy formulation in both combinatorial and physical contexts.

\paragraph{Why Does This Matter?}
Identifying blackbody radiation—a cornerstone of modern physics—as having an underlying structure equivalent to an NP-complete problem, and recognizing that nature effortlessly reaches equilibrium, presents a compelling question: \emph{How do physical systems effectively solve this computationally hard problem?}  As we will explore further, the transition from discrete combinatorics to continuous approximations at large \(N\), formalized by Rota’s theorem, combined with principles from information theory, hints at efficient pathways to navigate this complexity, potentially shedding light on the relationship between P and NP.


\subsection*{Shannon's Coding Theorem and the $O(n \log n)$ Solution}
\addcontentsline{toc}{subsection}{Shannon's Coding Theorem and the $O(n \log n)$ Solution}

Now that we have established the connection between physics and NP-Completeness through the lens of Rota's Entropy Theorem and the \textsc{Exact Cover} problem, we can bring in Shannon's Coding Theorem to show how a polynomial-time solution (\(O(n \log n)\)) might emerge. In essence, Shannon's theorem gives us a powerful guarantee that, for sufficiently large systems, there \emph{exists} a highly efficient encoding/decoding strategy whenever the entropy structure is known.

\subsection*{The Essence of Shannon's Coding Theorem}
\addcontentsline{toc}{subsection}{The Essence of Shannon's Coding Theorem}

\noindent
\textbf{Source Coding and Entropy.}  
Shannon's Coding Theorem, a cornerstone of information theory, addresses how to \emph{encode} (or compress) information produced by a probabilistic source. If a random variable \(X\) has outcomes \(\{x_1, x_2, \ldots, x_n\}\) with probabilities \(\{p_1, p_2, \ldots, p_n\}\) and entropy
\[
  H(X) \;=\; -\sum_{i=1}^{n}\,p_i \log_2\bigl(p_i\bigr),
\]
then Shannon's theorem tells us that we can compress the source \emph{down to} about \(H(X)\) bits per symbol (on average), with arbitrarily small error, given sufficiently large sample sequences.

Formally, we have:

\begin{theorem}[Shannon's Source Coding Theorem]
\label{thm:ShannonSource}
Let \(X\) be a discrete random variable with outcomes \(\{x_1, x_2, ..., x_n\}\) and associated probabilities \(\{p_1, p_2, ..., p_n\}\). Let \(H(X)\) be its Shannon entropy:
\[
  H(X) \;=\; -\sum_{i=1}^{n}\,p_i \,\log_2\bigl(p_i\bigr).
\]
For any \(\epsilon>0\) and \(\delta>0\), there exists a coding scheme for sequences of length \(N\) drawn from \(X\) such that:
\[
  H(X) \;\leq\; L \;\leq\; H(X) + \epsilon,
\]
where \(L\) is the average number of bits per outcome, and the probability of a decoding error is less than \(\delta\) when \(N\) is sufficiently large.
\end{theorem}

\noindent
\textbf{A Key Insight: Polynomial Time.}  
While Shannon's Coding Theorem is \emph{existential}—it does not prescribe the exact algorithm for constructing the perfect code—\emph{it does imply} that an efficient encoding/decoding strategy (e.g., with time complexity on the order of \(O(n \log n)\)) \textit{must exist} for large \(N\). Equally crucial is that typical practical coding/decoding algorithms (like Huffman, arithmetic, or more advanced block/source codes) \emph{do} run in polynomial time. Thus, if our physical system aligns with Shannon's framework (i.e., it satisfies Rota's axioms and thereby has a ``Shannon-like'' entropy), an \textbf{efficient} encoding or solution method must be available.

\subsection*{Applying Shannon's Coding to Rota's Entropy}
\addcontentsline{toc}{subsection}{Applying Shannon's Coding to Rota's Entropy}

Recall from earlier sections that Rota’s Entropy Theorem states:  
\[
  H(p_1,\dots,p_n) \;=\; C \,\sum_i p_i\,\ln \!\Bigl(\tfrac{1}{p_i}\Bigr),
\]
for some constant \(C\), whenever a system's “entropy function” satisfies his five axioms. This means that any valid physical entropy (such as the Boltzmann--Planck combinatorial count) is essentially a scaled version of Shannon’s entropy.

\paragraph{Connecting to a Discrete Random Variable.}  
Interpreting each \emph{microstate} of our physical system—like a configuration of quanta or a Sudoku arrangement—as a symbol produced by a discrete source \(X\), we obtain probabilities \(\{p_i\}\) for those symbols (microstates). The system’s entropy is then \(\sum p_i\ln(1/p_i)\) (up to a constant). By Shannon's Theorem:

\begin{itemize}
\item We can \emph{encode} typical microstate sequences with an average description length close to \(H(X)\).  
\item We can \emph{decode} or reconstruct these microstates from such codes, on average, in time linear or near-linear in the number of microstates, typically \(O(N\log N)\).  
\end{itemize}

\noindent
This is precisely what we need to move from a (seemingly) intractable NP-Complete counting problem to a polynomial-time approach. Instead of brute-forcing every microstate (which would be exponential), \emph{Shannon’s Coding Theorem guarantees} a compressed representation and efficient decoding in the limit of large \(N\).


\section*{\texorpdfstring{$P=NP$}{P=NP}!!! The Final Leap}
We established that the Blackbody Bus Station is equivalent to Sudoku and both are NP-Complete. We then showed that Rota’s Entropy Theorem provides a bridge between the Blackbody problem and \emph[all entropy problems], which means all these problems are NP-Complete AND vice versa. Finally, we showed that all these problems can be efficiently encoded and decoded using Shannon’s Coding Theorem. This leads us to the final conclusion: \(\boxed{P = NP}\)!!!.

By mapping a prototypical NP-Complete problem (e.g., \textsc{Exact Cover} in blackbody radiation) to an entropy count \(\log(W)\), \textbf{and} recognizing that Rota’s axiom-driven entropy must be of the Shannon form, we invoke Shannon’s Theorem to show an \emph{efficient} coding/decoding method. In purely computational complexity terms, this means \(\mathrm{NP} \subseteq \mathrm{P}\) at large scales—hinting strongly that \(P=NP\). The critical leap is the realization that the universe (and large combinatorial systems) \emph{are} governed by exactly these Shannon-like or Boltzmann--Planck distributions.

\paragraph{An Existence Proof, Not a Construction Manual.}  
As with many theorems of this nature, we do not yet have a step-by-step blueprint for building the universal \(O(n \log n)\) solver for \textsc{Exact Cover} (or for blackbody microstates). As, Scott Aaronson, one of the leading mathemticians in quantum computing says, "there is such a thing as a non-constructive proof ..."\cite{aaronson2021prove} and here that's exactly the role of Shannon’s Theorem - it \emph{ensures} \(P=NP\)'s existence. This parallels the way real physical processes—like thermalization—somehow “find” the stable entropy distributions without enumerating all possibilities. Similarly, the \emph{Coding Theorem} legitimizes a polynomial-time path to solutions of NP-Complete problems, even if we have not pinned down the explicit algorithm in every detail.

\section*{Discussion: Large-$N$ Physics and Being ``Forced to Do Good''}
\addcontentsline{toc}{section}{Discussion: Large-$N$ Physics and Being ``Forced to Do Good''}

So, does all of this mean we should drop everything and become champion Sudoku solvers or crypto hackers who can unlock Satoshi’s famous wallet? Not quite. Paradoxically, \emph{large-$N$} actually \textbf{protects} us when $n$ (or the system size) is too small to leverage the universe’s natural error-correction. This is part of why classical encryption isn’t trivially broken and why small-scale NP-Complete puzzles don’t instantly collapse. Yet it’s precisely in those \textbf{largest} problems—like protein folding, gravity-driven many-body systems, or truly massive prime factorizations—where nature and advanced AI \emph{do} display an uncanny efficiency. This section explores how large-$N$ ``forces us to do good,” bridging from the Principle of Least Action to scale invariance and practical coding methods like the FFT, and pondering how software and hardware might now scale together in unimaginable ways.

\subsection*{Principle of Least Action: Does Nature Solve Everything Efficiently?}
The \textbf{Principle of Least Action} has long hinted that physical systems seem to pick out optimal or near-optimal trajectories from infinitely many alternatives. From a discrete, NP-Complete perspective, it is astounding that planets, photons, and wavefunctions appear to “know” the best paths without enumerating all possibilities. Quantum mechanics goes a step further via path integrals, summing over exponentially many configurations yet yielding stable, predictable outcomes. This strongly resonates with our observation that, at large $N$, nature’s combinatorial processes \emph{must} converge efficiently to an extremal solution. Indeed, we might see the principle of least action as an embodiment of scale invariance: once you have enough quanta and interactions, “bad” solutions become vanishingly small in measure, and the system “selects” an optimal or stable microstate.

\subsection*{Forced to Do Good: Security at Small $n$, Efficiency at Large $N$}
One of the most gratifying aspects of this work is realizing that large-$N$ \emph{mandates} error at small $n$. In turn, this guarantees that problems like factoring a 2048-bit RSA number (relatively small in the cosmic sense) are still difficult in practice. This tension is actually a \emph{feature}: the universe’s computational complexity ensures our security, privacy, and even day-to-day stability. Yet when you scale up to truly massive $N$—as in astrophysics, plasma fusion, or large biomolecular systems—nature solves these NP-Complete problems in ways that appear polynomial time.

\subsection*{Recursive Partitioning and the Fast Fourier Transform}
\label{sec:FFTpartitioning}
\begin{quotation}\noindent
\emph{``Any Shannon entropy system can be partitioned with an FFT because of the deep mathematical connection of the Shannon entropy function and the discrete Fourier transform.''}
\end{quotation}

A concrete glimpse into efficient large-$N$ solutions is offered by \textbf{recursive partitioning algorithms}, such as the \textbf{Fast Fourier Transform (FFT)}. At first, the FFT might seem to have little to do with blackbody radiation or NP-Completeness. But recall our “bus station” analogy, where each bus frequency (``box'') is a prime-labeled mode, and each passenger (``ball'') is a discrete quanta of energy. Partitioning the occupancy distribution among these modes is mathematically analogous to \emph{decomposing} a large signal into its constituent frequencies. 

\begin{figure}[H]
    \centering
    \includegraphics[width=0.7\textwidth]{images/FFT_Partitioning.png}
    \caption{Recursive partitioning used by the FFT algorithm, illustrating how a signal is broken down into smaller segments for $O(n\log n)$ efficiency.}
    \label{fig:FFTPartitioning}
\end{figure}

As depicted in Figure~\ref{fig:FFTPartitioning}, the FFT recursively splits a signal into halves, computes partial solutions, and merges them. This strategy is exactly the “diamond in the sand” approach of dividing the search space into smaller parts until the solution emerges clearly. Crucially, the FFT has time complexity $O(n\log n)$—significantly faster than the naive $O(n^2)$ or $O(n^2\log n)$ methods of direct Fourier transforms. The analogy extends: once your “bus station” or partition problem is large enough, these splits average out errors so effectively that you approach a polynomial-time solution.

\subsection*{FFT as Shannon Coding for Prime Factorization}
Beyond decomposing signals, the FFT conceptually acts as a \textbf{Shannon coding} tool that separates mixed frequencies. Interpreted combinatorially, this is akin to prime factorization:

\begin{itemize}
\item \emph{Prime Labeling of Modes.} If each quantum mode is labeled by a prime $p_i$, then each occupancy pattern corresponds to a product of primes (like factoring an integer). 
\item \emph{Frequency Separation as Factor Extraction.} The FFT’s decomposition of overlapping signals matches the logic of extracting prime factors from a large number. Both tasks revolve around unraveling “hidden” base components.
\end{itemize}

Thus, the well-known $O(n\log n)$ efficiency of the FFT implies a similar bound on \emph{prime factorization} in a large-$N$, approximate sense—exactly matching our thesis that the universe’s combinatorial partitions can be solved in near-polynomial time for sufficiently large $N$. 

\subsection*{Why Public Key Encryption is (Mostly) Safe}
\label{sec:whyEncryptionSafe}
\paragraph{Errors at Small $n$ as Security.}
In practical cryptography, \emph{exact} factorization is mandatory to break RSA. Approximate solutions won’t do. The FFT and other partition-based methods remain approximate at modest sizes, so even if you’re “close” to finding the prime factors, you still fail to decrypt. For large $N$ in a physical sense (e.g., truly astronomical scales), factorization might become feasible, but that regime is out of reach for everyday cryptographic challenges.

\paragraph{Encryption as a Feature, Not a Bug.}
Hence, ironically, the same combinatorial partitioning that solves big cosmic-scale problems efficiently also protects our small-scale encryption by introducing persistent rounding or approximation errors. A factorization method off by even 1 in the prime guess is useless. So while $P=NP$ at the scale of the cosmos, cryptography remains robust for $n$ of a few thousand bits.

\subsection*{Scale Invariance in Action: From ATTRACTION to the Redwood Dilemma}
In our prior \textsc{ATTRACTION} paper \cite{AbadirAttraction2024}, we derived quantum gravity by harnessing \emph{scale invariance}, showing that at any scale, from subatomic to cosmic, the same combinatorial logic holds. There, we argued that continuous calculus fails (as shown by Poincaré for the three-body problem) yet a discrete combinatorial approach “just works.” 

\begin{itemize}
\item \textbf{Problems We Hoped REDWOOD Would Solve Are Precisely the Largest-$N$.} 
  Our hypothetical quantum computer \textsc{Redwood} was meant to tackle massive-scale NP-Complete challenges like protein folding, climate modeling, or general AI. But these are \emph{exactly} the cases where standard quantum computing’s continuum-based superpositions won’t help if $P \neq NP$. 
\item \textbf{But Large-$N$ \emph{Is} the Key.} 
  Once we accept discrete quanta and scale invariance, the problems that looked impossibly large ironically become the ones that nature (and perhaps advanced AI algorithms) can handle polynomially—just as the FFT partitions signals more efficiently the bigger and more structured they are. 
\end{itemize}

\subsection*{An Infinite Source of Bits: Energy Units, Partitions, and Hardware}
A breathtaking implication is that \emph{discreteness} itself might unlock a radically new synergy between \textbf{hardware} and \textbf{software}. If nature’s fundamental quanta come in endlessly divisible partitions (energy lumps that can keep subdividing), each partition potentially encodes bits—vastly expanding the horizon of computational hardware:

\begin{itemize}
\item \textbf{Software Gains from $P=NP$.} 
  Should we fully leverage the polynomial-time solutions guaranteed by Shannon coding, algorithms themselves might enjoy an exponential jump in effectiveness—no longer constrained by standard complexity barriers. 
\item \textbf{Hardware Gains from Discrete Scale Invariance.}
  At the same time, the discreteness and scale invariance of energy units suggest an \emph{infinite layering} of bits upon bits, fueling new hardware architectures. Imagine harnessing not superposition qubits but partition-based \textbf{energy bits} that scale with system size—offering a compounding synergy of hardware and software advancement.
\end{itemize}

In simpler terms, we might be on the cusp of a self-reinforcing cycle: as we discover better $O(n \log n)$ algorithms for large combinatorial tasks, we can also design hardware that subdivides energy quanta more effectively—thus \emph{increasing} our supply of discrete bits, which in turn accelerates algorithmic breakthroughs.

\subsection*{On the Possibility of Coding the Universe}
Finally, stepping back to the big picture: the results here do not imply that the universe is “completely solved.” Rather, they show a \emph{proof of existence} that \(\mathrm{Entropy} = P = NP\), meaning that large-scale real-world problems must be solvable in \(\mathcal{O}(n \log n)\) time. But we have not constructed a single, universal algorithm for everything. The Fast Fourier Transform, however, stands as a powerful metaphor (and partial reality) of how \emph{partition and conquer} is exactly what nature does. 

\begin{quotation}
\noindent
\emph{``Devising particular codes is a highly nontrivial task.''} \\
---\,Gian-Carlo Rota
\end{quotation}

That caution rings true. For “coding the universe,” we would need a systematic means of recursively subdividing not just signals or prime-labeled bus stations but \emph{all} quantum degrees of freedom—emergently producing physical laws, coupling constants, and boundary conditions. Such a vision remains speculative. Nonetheless, the scale-invariant, discrete combinatorial framework we have explored suggests a route to unify gravity, quantum mechanics, and complexity theory without resorting to artificial continuity or superposition. 

\subsection*{Conclusion: A Universe That Favors Efficiency at Scale}
\label{sec:discussionConclusion}
\begin{table}[H]
    \centering
      \begin{tabularx}{\textwidth}{|c|m{5cm}|m{3cm}|>{\centering\arraybackslash}m{2cm}|>{\centering\arraybackslash}m{2cm}|}
          \hline
          \textbf{Class} & \textbf{Description} & \textbf{Computational Complexity} & \textbf{Digital Computable} & \textbf{Quantum Computable} \\ \hline
          \textbf{P} & Solvable in polynomial time by a deterministic Turing machine (\(O(N^k)\)). & Polynomial (\(O(N^k)\)) & \checkmark & \checkmark \\ \hline
         
          \textbf{NP} & Nondeterministic Polynomial-Time; solutions verifiable in polynomial time. & Verifiable in \(\mathrm{poly}\) time & \checkmark (\(P=NP\)) & \checkmark (\(P=NP\)) \\ \hline
          \textbf{NP-Complete} & The hardest problems in NP; a polynomial-time solution to one solves all NP problems. & \emph{Would} be polynomial if \(P=NP\) & \checkmark (\(P=NP\)) & \checkmark (\(P=NP\)) \\ \hline
              
          \textbf{NP-Hard} & As hard as NP-Complete, not necessarily in NP. & At least as hard as NP-Complete &  &  \\ \hline
          \textbf{EXPTIME} & Problems growing exponentially vs.\ subdividable quanta. & Exponential (\(O(2^N)\)) &  & \checkmark?\\ \hline
      \end{tabularx}
      \caption{Revisiting the Hierarchy of Computational Complexity (Expanded for Energy-Unit Bits)}
      \label{tab:computability_expanded}
\end{table}

Table~\ref{tab:computability_expanded} incorporates our new insights. We see that physics is NP-Complete, yet effectively \emph{polynomial} at large $N$. The Redwood quantum supercomputer we once imagined becomes less relevant if $P=NP$: classical or discrete energy-partitioning approaches can already do the heavy lifting. Meanwhile, the scale invariance uncovered in \textsc{ATTRACTION} \cite{AbadirAttraction2024} and in partition-based methods like the FFT reveals a synergy between software algorithms ($P=NP$) and hardware-level discrete bits that could multiply our computing power beyond conventional speculation.

Returning to two guiding quotes:

\begin{quotation}
\textbf{Gian-Carlo Rota}:\quad
\emph{``What entropy does is to make the connection [between probability and information] precise... 
Finally, we turn to continuous random variables and prove that \emph{essentially all the interesting
distributions in probability theory} may be defined by entropy considerations.''}
\end{quotation}

\begin{quotation}
\textbf{Albert Einstein}:\quad
\emph{``The question seems to me to be how one can formulate statements about a discontinuum 
without resorting to a continuum... but for this we unfortunately are still lacking the mathematical form.''}
\end{quotation}

We have argued that Rota’s “entropy = counting” approach, bolstered by Shannon’s Theorem, provides precisely that missing “mathematical form” for the discontinuum Einstein sought. The deep partition logic at large $N$ explains how nature, AI, or future \emph{energy-bit} devices might out-compute naive superposition-based quantum hardware. In short, the \textbf{universe itself “codes” solutions to NP-Complete problems} when $N$ is big enough, verifying that \emph{yes}, $P=NP$ for the tasks that matter most at large scale—and leaving the small-scale realm to serve as our guardrail for cryptographic security and puzzle enjoyment.

\bigskip

\section*{Afterword: A Great Teacher And Giver of Tricks}
\addcontentsline{toc}{section}{Afterword: A Great Teacher And Giver of Tricks}
In memory of the master who not only me my one trick to play against all problem, but, more importantly taught with an infectious and eye-opening enthusiasm. I'm truly sorry it took me 30 years to see what you put in front of me.
\begin{quotation}
    \emph{A teacher can be effective in discovering his students' talents and in encouraging those talents. A good teacher does not teach facts; he teaches enthusiasm, open-mindedness, and values.}

    ...

    \emph{All creative mathematicians and physicists I know work like this. They have constantly in their minds a list of a dozen of their pet problems. Anything they hear, they automatically test against their dozen problems. Every once in a while, there's a hit, and then people cry out, "He's a genius! How did he know that this was the right trick to apply to Problem Three?" Little do they know that the guy has been testing for years everything he hears against Problem Three. That's the way a lot of discoveries are made.}
    
    ...

    \emph{Research is sometimes not so much discovering something new as becoming aware of prejudices that stop us from seeing what is in front of us.}
    \newline
    \newline
    
    - \textbf{Gian-Carlo Rota}, \textit{Mathematics, Philosophy, and Artificial Intelligence: a dialogue with Gian-Carlo Rota and David Sharp}, 1985\cite{Rota1985} (emphasis added) 
\end{quotation}

\part{Addendum: Supplementary Mathematical Proofs \& Physics Derivations}

\section{Rota's Entropy Theorem: Proof and Implications}

\paragraph{Proof of Rota's Entropy Theorem}
The proof of Rota's Entropy Theorem is a key result in the theory of information and entropy, and it is an honor to share it here.

The remainder of this section is excerpted from class text provided by Professor Gian-Carlo Rota. To my knowledge it is unpublished and uncopyrighted. "Introduction to Probability Theory, Second Preliminary Edition" manuscript circa 1993, authors are Kenneth Baclawski, Gian-Carlo Rota, \& Sara Billis. It is similar to the one on the Internet Archive\cite{RotaBaclawski} where the same proof is present, but I have not found this particular version online. 
\medskip
\medskip

\subsection*{Chapter VIII: Entropy and Information}
\normalsize
\subsection*{Properties of Entropy}

So far, we have discussed examples of the entropy of some random variables. Although these examples provide some motivation for our definition of entropy, they leave unanswered the more difficult question of why, out of all possible definitions, we use this one.

We will do this by finding five self-evident properties that ought to hold for any reasonable measure of information (or entropy). It then turns out that our definition of entropy is the only one that satisfies all these properties.

We begin with the most obvious of properties. As we have defined it, \( H \) is a function of partitions of the sample space. However, it should be clear that we want \( H \) to depend only on the set of probabilities of the blocks of the partition. In fact, we want \( H \) to depend only on the positive probabilities which occur. Moreover, we want \( H \) to be a continuous function of these probabilities. This is a convenience only. We could, with a great deal of effort, derive continuity from other more complex conditions; but we would rather concentrate on the important issues. 

We summarize the conditions on \( H \) we have just described before going on to the difficult question of conditional entropy.

\paragraph{Entropy Property 1:} 
An entropy is a function defined on sets \( \{p_1, p_2, \ldots, p_n\} \) of non-negative real numbers, which satisfy \( p_1 + p_2 + \ldots + p_n = 1 \).

\paragraph{Entropy Property 2:} 
If \( H \) is an entropy function, then for any set \( \{p_1, p_2, \ldots, p_n, 0\} \) on which \( H \) is defined, \( H \) satisfies:
\[
H(p_1, p_2, \ldots, p_n, 0) = H(p_1, p_2, \ldots, p_n).
\]
In other words, \( H \) depends only on the nonzero \( p_i \)'s in a given set.

\paragraph{Entropy Property 3:} 
An entropy function is continuous.
The next property of entropy we consider requires the concept of conditional entropy. There are two ways to think of conditional entropy, and the fact that they are equivalent is our next property of entropy. To illustrate the ideas involved, we consider the following simple weighing problem:

We have three coins, some of which may be counterfeit (but not all). Counterfeit coins are distinguishable from normal coins by the fact that they are lighter. We are given a balance scale, and we wish to find out which, if any, of the coins are counterfeit. The sample space for this problem consists of seven sample points, one for each possible set of good coins. We denote them as follows:
\[
\Omega = \{1, 2, 3, 12, 13, 23, 123\}.
\]

Now what happens when we put the first two coins on each side of the scale? The sample space is partitioned into three blocks corresponding to the three possible outcomes of the weighing:
\[
\sigma = \{\{12, 123, 3\}, \{2, 23\}, \{1, 13\}\}.
\]
After recording the result of this weighing, we then place the second and third coins on the two sides of the scale. The result of this second weighing is to partition each of the blocks of the first weighing:
\[
\{12, 123, 3\} \to \{\{12\}, \{123\}, \{3\}\}, \quad 
\{2, 23\} \to \{\{2\}, \{23\}\}, \quad 
\{1, 13\} \to \{\{1\}, \{13\}\}.
\]

The combined information of the two weighings is represented by the partition into seven blocks, each with one sample point. Call this partition \( \pi \). Conditional entropy is concerned with the effect of the second weighing, given that the first has occurred. One way to analyze this is to look at each block \( \sigma_i \) of the partition of the first weighing and to analyze the situation as if it were the whole sample space. In general, for an event \( A \) and a partition \( \tau \), we define the conditional entropy of \( \pi \) given \( A \), written \( H(\pi|A) \), to be the entropy of the partition \( \tau_1 \cap A, \tau_2 \cap A, \ldots \) that \( \tau \) induces on \( A \). 

Thus, in the above weighing problem, we have three conditional entropies, one for each possible outcome of the first weighing:
\[
H(\pi|\sigma_1), \quad H(\pi|\sigma_2), \quad H(\pi|\sigma_3).
\]
The conditional entropy of \( \pi \) given \( \sigma \) is then defined to be the average of these. More precisely, if \( \pi \) and \( \sigma \) are any two partitions of a sample space \( \Omega \) such that \( \pi \) is finer than \( \sigma \), we define the conditional entropy of \( \pi \) given \( \sigma \) to be the average value of \( H(\pi|\sigma_i) \) over all blocks \( \sigma_i \) of \( \sigma \):
\[
H(\pi|\sigma) = \sum_i P(\sigma_i) H(\pi|\sigma_i).
\]

On the other hand, we would like to think of information as a "quantity" that increases as we ask more and more questions about our experiment. Therefore, the conditional entropy of \( \pi \) given \( \sigma \) ought to be the net increase in entropy from \( \sigma \) to \( \pi \). In other words, we require our entropy function to satisfy:

\paragraph{Entropy Property 4:} 
If \( \pi \) is a finer partition than \( \sigma \), then
\[
H(\pi|\sigma) = H(\pi) - H(\sigma).
\]

The last property we require is one that we have already discussed. The partition having maximum entropy among all partitions with a given number of blocks is the one for which all the blocks have the same probability.

\paragraph{Entropy Property 5:} 
If \( H \) is an entropy function, then any set \( \{p_1, p_2, \ldots, p_n\} \) on which \( H \) is defined satisfies:
\[
H(p_1, p_2, \ldots, p_n) \leq H\left(\frac{1}{n}, \frac{1}{n}, \ldots, \frac{1}{n}\right).
\]

We are now ready for the following remarkable fact: if \( H \) satisfies the above five properties, then \( H \) is given by the formula introduced earlier in this chapter, except for a possible scale change.

\subsection*{Uniqueness of Entropy}
If \( H \) is a function satisfying the five properties of an entropy function, then there is a constant \( C \) such that \( H \) is given by:
\[
H(p_1, p_2, \ldots, p_n) = C \sum_{i} p_i \log_2 \frac{1}{p_i}.
\]

\paragraph{Proof:} 
The proof is rather technical, so we suggest omitting it on the first reading. However, it is of interest to outline the main points. To show that \( H \) has the form given above, we use the following two facts:

1. The entropy of the partition consisting of just one block of probability \( 1 \) is zero, i.e., \( H(\Omega) = 0 \). By definition, \( H(\Omega) \) is the same as \( H(\{1\}) \). Therefore, \( H(\Omega) = H(\{1\}) = 0 \).

2. We define a function \( f(n) \) by \( H\left(\frac{1}{n}, \frac{1}{n}, \ldots, \frac{1}{n}\right) \). We have just shown that \( f(1) = 0 \) and we want to calculate \( f(n) \) in general. Using properties 2 and 5, we show that \( f(n) \) is increasing:
\[
f(n) \leq f(n+1).
\]

Next, we consider a partition \( \sigma \) consisting of \( n^k \) blocks, each of which has probability \( \frac{1}{n^k} \). Then subdivide each of these into \( n \) parts, each of which has the same probability. Call the resulting partition \( \pi \). The conditional entropy \( H(\pi|\sigma) \) for each block \( \sigma_i \) is clearly given by \( f(n) \). Thus the conditional entropy \( H(\pi|\sigma) \) is \( f(k) - f(k-1) \). If we apply this fact \( k \) times, we obtain:
\[
f(n^k) = k f(n).
\]

Now fix two positive integers \( n \) and \( k \). Since the exponential function is an increasing function, there is an integer \( b \) such that:
\[
2b \leq n^k < 2b+1.
\]
We now apply the two facts about \( f(n) \) obtained above to this relation:
\[
f(2^b) \leq f(n^k) \leq f(2^{b+1}).
\]
Since \( f(n) \) is increasing, we know:
\[
b f(2) \leq k f(n) \leq (b+1)f(2).
\]
Now divide these inequalities by \( k f(2) \):
\[
\frac{b}{k} \leq \frac{f(n)}{f(2)} \leq \frac{b+1}{k}.
\]
Now apply the increasing function \( \log_2 \) to the inequalities:
\[
\frac{b}{k} \leq \log_2(n) \leq \frac{b+1}{k}.
\]
It follows that both \( f(n)/f(2) \) and \( \log_2(n) \) are in the interval \( [b/k, (b+1)/k] \). This implies that \( f(n)/f(2) \) and \( \log_2(n) \) can be no farther apart than \( 1/k \), the length of this interval. But \( n \) and \( k \) were arbitrary positive integers. So if we let \( k \) get very large, we are forced to conclude that:
\[
f(n)/f(2) = \log_2(n).
\]
Thus, for positive integers \( n \), we have:
\[
f(n) = f(2) \log_2(n).
\]

We will define the constant \( C \) to be \( -f(2) \). Since \( f(2) \geq f(1) = 0 \), we know that \( C \) is negative.

We next consider a set \( \{p_1, p_2, \ldots, p_n\} \) of positive rational numbers such that \( p_1 + p_2 + \ldots + p_n = 1 \). Let \( N \) be their common denominator, i.e., \( p_i = \frac{a_i}{N} \) for all \( i \), where each \( a_i \) is an integer and \( a_1 + a_2 + \ldots + a_n = N \). Let \( \sigma \) be a partition corresponding to the set of probabilities \( \{p_1, p_2, \ldots, p_n\} \). Let \( \pi \) be a partition obtained by breaking up the \( i \)-th block of \( \sigma \) into \( a_i \) parts. Then every block of \( \pi \) has probability \( \frac{1}{N} \). By definition of conditional entropy:
\[
H(\pi|\sigma) = -\sum_{i} P(\sigma_i) H(\pi|\sigma_i) = -\sum_{i} f(a_i) - C \sum_{i} p_i \log_2(a_i).
\]

By property 4, on the other hand, we have:
\[
H(\pi|\sigma) = H(\pi) - H(\sigma) = f(N) - H(\sigma).
\]
Combining the two expressions for \( H(\pi|\sigma) \) gives us:
\[
H(\sigma) = -C \log_2(N) + C \sum_{i} p_i \log_2(a_i).
\]
By continuity (property 3), \( H \) must have this same formula for all sets \( \{p_1, p_2, \ldots, p_n\} \) on which it is defined. This completes the proof.

We leave it as an exercise to show that the above formula for entropy actually satisfies the five postulated properties. We conclude by giving an interpretation of independence of partitions in terms of conditional entropy. Intuitively, if \( \pi \) and \( \sigma \) are independent, then their joint entropy \( H(\pi \cap \sigma) \) is the sum of the individual entropies:
\[
H(\pi \cap \sigma) = H(\pi) + H(\sigma).
\]
In terms of conditional entropy, this says that \( H(\pi \cap \sigma) = H(\pi) \).

\subsection*{The Shannon Coding Theorem}

A consequence of Entropy Property 4 of the last section is that if we wish to answer a question \( X \) by means of a sequence of questions \( S_1, S_2, \ldots, S_n \), the joint entropy of \( S_1, S_2, \ldots, S_n \) must be at least as large as the entropy of \( X \), and hence the sum of the entropies of the \( S_i \)'s must be at least as large as the entropy of \( X \). In particular, if the \( S_i \)'s are yes-no questions, then \( H_2(S_i) \leq 1 \) and we get the crude inequality:
\[
n \geq H_2(X).
\]

The problem of finding a set of sufficient statistics for a random variable \( X \) is called the \emph{coding problem} for \( X \), and the sequence \( S_1, S_2, \ldots, S_n \) is said to \emph{code} \( X \). As we will see in the exercises, the kinds of questions one may ask are usually restricted to some class of questions. Devising particular codes is a highly nontrivial task.

One of the reasons that coding is so nontrivial in general is that one is usually required to answer a whole sequence of questions \( X_1, X_2, \ldots \), produced by some process, and as a result one would like to answer the questions in the most efficient way possible. Consider one example. Suppose that \( X \) takes values \( 1 \) through \( 200 \) each with probability \( 0.85 \), and takes values \( 0 \) with probability \( 7.5 \times 10^{-4} \). Then \( H_2(X) \) is less than \( 1 \). Simply by counting one can see that at least \( 8 \) yes-no questions will be needed to achieve a sufficient statistic for \( X \), even though the entropy suggests that one should be able to determine \( X \) with a single yes-no question.

\section*{Quantum Gravity: Resolution of The EPR Paradox \& The Double Slit Experiment}
\addcontentsline{toc}{section}{Quantum Gravity: Resolution of The EPR Paradox \& The Double Slit Experiment}

\subsubsection*{NB: On the Linearity of Paths vs. Linear "Betting Odds" in the Double Slit Experiment}
\addcontentsline{toc}{subsubsection}{NB: On the Linearity of Paths vs. Linear "Betting Odds" in the Double Slit Experiment}
Previously in \textsc{Attraction} we spoke of linearity with respect to particles moving in a straight line like the bullets in Feynman's example-this was an oversimplification of linearity. A full mathematical treatment of linearity as it relates to the EPR paradox is not strictly needed for our proof below but it may still be intuitively helpful to discuss.  While it is beyond the scope of this paper to delve into the probabilities of gambling it may be interesting to the reader to peruse the introduction to Rota's Introduction to Probability Theory\cite{Rota1997} where he discusses Coin Tossing and the expactations of random walks - some very unexpected results occur from simple, fair coin tosses.\cite{Rota1985} 

Einstein, Podolsky, and Rosen (EPR) challenged the completeness of Quantum Mechanics because they thought that unexpected experimental results in the DSE must be the result of an incomplete understanding of physics, and not by some violation of physics as we know it. We can think of linear probabilities as the types of betting odds that happen at a casino or a horse track-people shouldn't be allowed to "beat the odds." In response to the EPR paradox, John Stewart Bell formulated his now-famous "Bell's Theorem", which derives a set of inequalities that must hold for any theory with linear correlations (i.e. "fair" betting odds). Imagine a horse-race betting site that sets "fair" odds for each horse based on standard assumptions. In this scenario, there are three main ways someone could cheat so that the final outcome doesn't match those odds. 
\begin{enumerate}
    \item \textbf{Non-locality:} the bettor receives the race result faster than the house's update (as though using a faster than light telephone signal) — analogous to “spooky action at a distance.”
    \item \textbf{Hidden variables:} the bettor starts the race with secret information the house doesn’t possess — like predetermined initial states for particles.
    \item \textbf{Reality isn’t fixed:} the bettor has a conspirator at the track who effectively \textit{decides} the race outcome in the bettor’s favor—meaning the “real” winner doesn’t exist independently until it’s reported.
\end{enumerate}
Bell’s inequality says that if your system is \textit{local} (no faster-than-light signals), \textit{realist} (the outcome exists in advance), and has \textit{linear} correlations, then you can’t systematically beat the odds. In essence, Bell formalized that if those assumptions are true then all relationships must be linear (with no extra "cheats").

However, Bell’s analysis assumes a \textit{linear} relationship to begin with, and it is that assumption that we choose to challenge in this paper. If the gambling system allows for \textit{non-linear} correlations—akin to reshaping how odds are calculated—then you could exceed those classic odds limits without resorting to any of the above “cheats.” In other words, you’d look like you were gaming the system under a linear model, but you wouldn’t actually be breaking locality or realism; you’d simply be using a different (non-linear) probability rule that Bell’s inequality does not address. Specifically, if your model has three properties--non-linearity, locality, and realism--then it is able to explain the Double Slit experiment without invoking concepts such as non-locality and superposition, and hence, we might conclude that Copenhagen based Quantum Mechanics is incomplete. 

Our logic proof below is premised on known behaviors of strange attractors in non-linear dynamic systems. While we believe non-linearity is the mechanism that lets particles "beat the odds," it might be that some other mechanism is at play and including this argument would unnecessarily complicate the proof.

\subsection*{Formal Proof: Quantum Gravity Completes Quantum Mechanics \& Resolves EPR Paradox}
\addcontentsline{toc}{subsection}{Formal Proof: Quantum Gravity Completes Quantum Mechanics \& Resolves EPR Paradox}
\textbf{Definitions:}
\begin{enumerate}
    \item \emph{Realism}: Physical properties exist independently of observation. The same rules apply across all scales.
    \item \emph{Scale Invariance}: The form of physical laws is invariant under changes in scale.
    \item \emph{No Local Hidden Variables}: No pre-encoded instructions or carried memory determine measurement outcomes.
    \item \emph{Locality (No Signaling)}: No information travels faster than light.
    \item \emph{Chaos and Strange Attractors}: Deterministic dynamics with sensitive dependence on initial conditions lead to self-organized, stable patterns without predefined states or signaling.
    \item \emph{Phenomena to Explain}: 
    \begin{itemize}
        \item $DSE$: The Double Slit Experiment pattern.
        \item $BELL$: Bell-type correlations.
    \end{itemize}
    \item $S(x)$: “Superposition/entanglement is strictly necessary for $x$.”
    \item $P(x)$: “A realistic, scale-invariant, chaotic, no-hidden-variable, local model (including quantum gravity) explains $x$.”
    \item $CI$: Copenhagen Interpretation.
    \item $QG$: Quantum Gravity (a scale-invariant gravitational interaction at quantum scales).
    \item $E(m)$: “Model $m$ is an incomplete understanding of Quantum Mechanics.”
\end{enumerate}

\textbf{Axioms:}
\begin{enumerate}
    \item \textbf{Realism and Scale Invariance Axiom}: The same physical principles (e.g., gravity, chaos) apply at all scales.
    \item \textbf{Locality Axiom}: $\forall$ spatially separated systems $(A_i,A_j)$, no superluminal communication influences outcomes.
    \item \textbf{No Shared Initial State Axiom}: Systems do not begin with pre-coordinated instructions determining all future measurements.
    \item \textbf{Chaos Axiom (3+ Body Quantum Gravity)}: With three or more particles interacting gravitationally at the quantum scale, deterministic rules plus sensitive dependence on initial conditions produce complex, fractal patterns (strange attractors).
\end{enumerate}

\subsubsection*{Proof By Contradiction}

\textbf{Goal}: To show that assuming superposition/entanglement as necessary leads to a contradiction if a chaotic, no-hidden-variable model with quantum gravity exists.

\textbf{Step 1 (Copenhagen Completeness Claim)}:
 Initially, the Copenhagen Interpretation (CI) is taken as complete if and only if superposition and/or entanglement are strictly necessary to explain both the Double Slit Experiment (DSE) and Bell-type correlations (BELL). Formally:
\[
S(DSE) \wedge S(BELL) \iff \neg E(CI),
\]
where $\neg E(CI)$ denotes that the CI is not incomplete (i.e., it is complete). This ensures that if we find a scenario where DSE and BELL can be explained without strict superposition/entanglement, CI cannot be considered complete.

\textbf{Step 2 (Alternative Model with Quantum Gravity)}:
Introduce a model $M = CI + QG$ that incorporates quantum gravity at small scales and is:
\begin{itemize}
    \item Realistic and Scale-Invariant.
    \item Local (no signaling).
    \item Without hidden variables (no pre-encoded outcomes).
    \item Chaotic: With three or more particles under quantum gravity, deterministic but highly sensitive initial conditions produce complex, fractal patterns (strange attractors).
\end{itemize}

If $M$ can reproduce the DSE and BELL phenomena without appealing to superposition/entanglement, we have:
\[
P(DSE): M \text{ explains DSE without } S(DSE),
\]
\[
P(BELL): M \text{ explains BELL without } S(BELL).
\]

\textbf{Step 3 (Chaos Implies No Local Hidden Variables)}:
By the Chaos Axiom, complex correlated outcomes emerge from deterministic chaos and gravitational interactions without any pre-encoded information. Thus, no local hidden variables are necessary. In other words, chaotic gravitational dynamics can yield complex, correlated patterns that mimic the statistical signatures of quantum phenomena without assuming any initial encoding of measurement outcomes.

\textbf{Step 4 (Implications of $P(x)$)}:
If $P(DSE)$ holds, we have:
\[
P(DSE) \implies \neg S(DSE).
\]
Similarly, if $P(BELL)$ holds:
\[
P(BELL) \implies \neg S(BELL).
\]

Combining both:
\[
P(DSE) \wedge P(BELL) \implies \neg S(DSE) \wedge \neg S(BELL).
\]

\textbf{Step 5 (Contradiction)}:
Recall from Step 1:
\[
S(DSE) \wedge S(BELL) \iff \neg E(CI).
\]
If $M$ exists and explains both DSE and BELL without superposition/entanglement, then:
\[
\neg S(DSE) \wedge \neg S(BELL).
\]

This directly contradicts the initial assumption:
\[
S(DSE) \wedge S(BELL).
\]

Since $S(DSE) \wedge S(BELL)$ was equivalent to $\neg E(CI)$, the contradiction shows that we cannot maintain that CI is complete if a $QG$ model explains the phenomena without strict superposition/entanglement. Hence, the initial assumption that CI is complete (dependent on superposition/entanglement being necessary) must be false.

\subsection*{Discussion: Intuition of Quantum Gravity Completing Quantum Mechanics}
\addcontentsline{toc}{subsection}{Discussion: Intuition of Quantum Gravity Completing Quantum Mechanics}

By positing a realistic, scale-invariant, local, chaotic model with quantum gravity that explains DSE and BELL correlations without superposition or entanglement, we have derived a logical contradiction with the assumption of CI's completeness. This supports the notion that CI may be incomplete. Rather than invoking superposition and entanglement, the necessary complexity and correlations can emerge naturally from chaotic gravitational dynamics. In subsequent sections, we will illustrate how this framework can also be used to derive Planck’s Law under quantum gravity, further strengthening the argument that Einstein’s intuition about a more complete, scale-invariant understanding of quantum phenomena may have been correct.

\section*{Derivation: Planck's Law Under Quantum Gravity \& Chaos}
\addcontentsline{toc}{section}{Derivation: Planck's Law Under Quantum Gravity}
Having shown that quantum gravity, chaos and scale-invariance can resolve the EPR paradox, a question naturally arises: can this new framework reproduce the foundational laws of physics? If the same logic applies across all scales, then it must also yield results that are consistent with what has been experimentally confirmed. We now turn our attention to a notoriously difficult problem in physics—the derivation of Planck's Law, the very equation that marked the birth of quantum theory. But where do we even begin? Perhaps by looking at what has already been ruled out and understanding how Planck himself arrived at his breakthrough.

\subsection*{Failed Correspondence, Or, "Calculus: The Wrong Tool For The Job"}
\addcontentsline{toc}{subsection}{Failed Correspondence, Or, "Calculus: The Wrong Tool For The Job"}
As we rigorously explored in the FRAQTL paper, calculus is an absolute impediment to describing a physics that includes Chaos Theory\cite{FRAQTL2024} since the very field of chaos was the result of Poincar\'e showing that gravitational N-Body problems have no closed-form (calculus) solution. But if calculus isn't the right tool, then what is? To find a different path forward, and to further reinforce that calculus is not the right tool for this problem, it's worthwhile to retrace some of the unsuccessful efforts of the 1980s and 90s when physicists attempted to apply the Copenhagen Interpretation to classical physics.

The "correspondence principle" of Niels Bohr, not surprisingly, states that Quantum Mechanics explains all other physics. Recall that the explicit assumption of Copenhagen is that the quantum realm is continuous and that calculus is its mathematics. In the 1980s and 90s, around the time Chaos Theory was gaining prominence, a related field called Quantum Chaos never quite took off. Michael Berry, one of the central figures in Quantum Chaos, penned a paper sub-titled "Is the Moon There When Somebody Looks?" and explained: "According to the correspondence principle, the classical world should emerge from the quantum world whenever Planck’s constant $h$ is negligible." Quantum Chaos was essentially an attempt to use the Copenhagen Interpretation to describe classical physics where gravity-induced chaos reigns (the N-Body problem). Instead of validating the correspondence principle, it served to prove the opposite; that the Copenhagen based Quantum Mechanics is fundamentally flawed \cite{Zurek1998}. The falsification of Copenhagen is semi-popularly known as "the problem with Hyperion" (a moon of Saturn). Noted quantum physicist, Sean Carroll, gives one of his typically elegant overviews:

\begin{quotation}
    "[Correspondence] is one of those vague but invaluable rules of thumb that was formulated by Niels Bohr back in the salad days of quantum mechanics. If it sounds a little hand-wavy, that’s because it is.
    \newline
    \newline
    The vagueness of the correspondence principle prods a careful physicist into formulating a more precise version, or perhaps coming up with counterexamples. And indeed, counterexamples exist: namely, when the classical predictions for the system in question are chaotic. \dots
    \newline
    \newline
    Some years ago, Wojciech Zurek and Juan Pablo Paz described a particularly interesting real-world example of such a system: Hyperion ... The orbit of Hyperion around Saturn is fairly predictable; happily, even for lumpy moons, the center of mass follows a smooth path. But the orientation of Hyperion, it turns out, is chaotic \dots
    \newline
    \newline
    So, in the real world, not only does this particular moon (of Saturn) exist when we’re not looking, it’s also in a pretty well-defined orientation \dots
    "\cite{CarrollHyperion}
\end{quotation}

Our objective is to find a mathematical framework for Quantum Mechanics upon which we can derive Planck's Law and other foundational relationships. As we just saw, the opposite could not be done—Quantum Chaos could not derive classical mechanics using Bohr's correspondence principle, the Copenhagen Interpretation, and calculus. This failure is a direct consequence of the foundation of Chaos Theory as formulated by Poincar\'e: chaos appears if gravity present, and chaos cannot be described by calculus. So, if calculus cannot describe systems where gravity appears, then how can it be the basis for a complete theory of physics?

It might seem that Quantum Chaos ruled out the possibility of Quantum Mechanics and classical physics from ever playing nicely together. However, in fact, all it did was rule out the Copenhagen Interpretation and its continuous assumptions. That is, if applying the Copenhagen Interpretation to classical physics fails, then the Copenhagen Interpretation is incomplete. This is the same conclusion we reached in the proof above and that proof was premised on a kind of "reverse-correspondence": that Chaos Theory can be applied to Quantum Mechanics and that doing so would make Quantum Mechanics complete.

\subsection*{Combinatorics: A Mathematical Form For A Discontinuum}
\addcontentsline{toc}{subsection}{Combinatorics \& The Three-Body Problem}
The Three-Body problem is a mythical beast of ornery mathematics, a problem that brought the great Newton himself to despair. With the idea of "reverse-correspondence" in mind, we opened a new door thinking it promising, but now we face the Three-Body monster head-on. If the goal is a theory of quantum gravity, should we try a "Chaotic Quantum Theory" of reverse-correspondence? The very fact that the Three-Body problem resists calculus should be a clear signal that we need a new mathematics, a mathematics capable of handling the discrete, the chaotic, and the discontinuous.

Letting history be our guide once more, we'll explore whether Planck's own derivation might give us a clue that will put this beast to sleep. Michael Fowler's wonderful recounting of Planck's work reveals: \cite{Fowler2008}

\begin{quotation}
    Planck was not impressed by [Boltzmann's] approach, since it implied that the Second Law was only statistical, only valid in the limit of large systems. However, he was fairly familiar with Boltzmann’s work — he’d taught it in some of his classes, to present all points of view. Planck wasn’t sure, though, that atoms and molecules even existed.

    Planck did this work by December 1900, in two intense months after learning the new experimental results and feeling he had to justify his curve that fit so well. But he only half believed it. After all, the first part of his derivation, identifying the energy of an oscillator with that in the radiation field, was purely classical: \textit{he’d assumed the emission and absorption of energy to be continuous}. Then, he suddenly changed the story, moving to a totally nonclassical concept, that the oscillators could only gain and lose energy in chunks, or quanta. (Incidentally, it didn’t occur to him that the radiation \textit{itself} might be in quanta: \textit{he saw this quantization purely as a property of the wall oscillators}.) As a result, although the exactness of his curve was widely admired, and it was the Birth of the Quantum Theory (with hindsight), no-one — including Planck — grasped this for several years!
\end{quotation}

After nearly a century of debate and attempts to unify the quantum and classical regimes, finding a theory of quantum gravity remains a formidable challenge. With that in mind, we will do as Planck did: we will use Boltzmann's methods, which, like all of combinatorics, is the mathematics of counting discrete possibilities, and we will blithely ignore the Three-Body beast... maybe it will simply go away. The approach here is to not simply "add more math" to a problem, but to ask what math is most appropriate for the job. If calculus, the math of the continuum, has failed, then what other mathematical tool should we use for a discontinuum?

The answer, it turns out, was already waiting for us in the work of Boltzmann—the mathematics of combinatorics. Combinatorics is fundamentally about counting the number of ways things can be arranged, the number of possibilities in a system. It is the mathematics of "balls-into-boxes", a discrete framework that naturally fits with a vision of a universe that is not continuous but granular, a universe made up of individual, countable quanta. While calculus deals with smooth curves and continuous change, combinatorics is perfectly suited to handle discrete objects and their interactions, and as such, it also suits our understanding of chaos.

By making the switch to a combinatorial perspective, we will bypass the difficulties inherent in the Three-Body problem, allowing us to proceed in a new, more tractable direction. Instead of attempting to solve complex differential equations, we will focus on counting the ways that energy can be distributed, a method that mirrors Planck's own approach when he first introduced the idea of quantization. Instead of asking what *can* happen, we will look at what is *most probable* to happen, just as Boltzmann did when he derived his famous H-Theorem. This move to combinatorics is not a random choice; it is a deliberate shift in our mathematical toolkit to match the discrete and discontinuous nature of reality itself, it is a move that we argue is necessary to bring gravity into quantum mechanics.

\subsubsection*{A Change In Perspective: Discrete \& Divisible Quanta:}
\addcontentsline{toc}{subsection}{A Change In Perspective: Discrete \& Divisible Quanta}
To recap, our objective is to bring gravity into Quantum Mechanics because we know that (Copenhagen based) Quantum Mechanics cannot be brought into classical gravity. In our previous proof we showed that chaos can resolve the EPR paradox, but we also know that the N-Body problem underlying chaos is incredibly difficult. To proceed, we are taking a page out of Planck's derivation of his own law, proceeding with Boltzmann's combinatorics to get over the Three-Body problem, and attempting a new approach to Quantum Mechanics that relies on discreteness.

Since combinatorics is the math of "balls-into-boxes," we have a new problem: what exactly are the balls, and what are the boxes? In Copenhagen Quantum Mechanics, particles are treated as point-like with no size and potentially occupying multiple places at once. Therefore, before we begin, we need to re-examine the foundational continuum assumption of the Copenhagen Interpretation and, once again, follow Einstein's intuition.

\begin{quotation}
    "Electrons (as points) would be final conditions (building blocks) in such a system. Do final building blocks exist at all?  Why are they all similar in size? Is it adequate to say: God in His wisdom has made them all the same size, each like the others, because He so wished? If it had suited Him, He could also have made them dissimilar." \cite{Einstein314}
\end{quotation}

In a discontinuum, every particle—like photons and electrons—must have a finite size and a definite position. Instead of treating photons as point-like (having no size), or as wave which is everywhere all at once, we will treat them each as a "box" having definite size and position. In fact, we'll even treat the quanta as divisible into smaller "energy unit" components that have no name in current physics (each a "ball"), with the understanding that these smaller components also possess the same properties. This is a departure from the current thinking, yet also a necessary step to get to a useful result.

Now, having changed our perspective we can add Einstein's 1905 toolkit, including quanta (discrete units which are divisible), Brownian motion (stochastic oscillators), Relativity frames, and the mass-energy relation $E=mc^2$. These subdividable quanta (frames) under stochastic motion and an attractive force will be shown to yield Planck’s Law.


\subsection*{A Blackbody Thought Experiment: The Blackbody Bus Station}
\addcontentsline{toc}{subsection}{A Blackbody Thought Experiment: The Blackbody Bus Station}

The discussion up to this point has seemed rather removed from an intuitive understanding of what is happening at the quantum level. Without such an intuitive understanding balls-into-boxes problems get very painful very quickly. Taking another page out of Einstein's book, let's try an accessible analogy. While Einstein often used a railway station as a conceptual tool, here we will consider a bustling `Blackbody Bus Station.'' Our aim is to understand the relationship between the number of passengers in the station and the fullness of buses when they depart. In this analogy, buses are colored according to how full they are: emptier buses appear red, and as they become more crowded, they transition through the spectrum towards blue.

At its core, the blackbody problem is about predicting the most probable colors of light emitted by a hot object. In more technical terms, color corresponds to the energy distribution of emitted light, also known as its spectral intensity. If an object is heated to 1000 K, will it emit more red light or more blue light? If it is heated further to 2000 K, how does the relative intensity of emitted red versus blue light change? Planck’s Law provides the quantitative relationship between temperature and spectral intensity.

To connect these ideas to our thought experiment, we will make parenthetical notes that relate station elements to their physical and mathematical counterparts. In the Blackbody Bus Station, the passengers (energy units) and bus drivers have been partying a bit too hard. As a result, they stagger about randomly, reflecting a stochastic random walk (Brownian motion). The passengers hop on and off buses (quanta) in seemingly chaotic ways, and the bus drivers, equally disoriented, swerve unpredictably (another manifestation of Brownian motion). All the while, the drivers shout at their passengers to move towards the center of the bus (an attractive force). A strict speed limit applies everywhere (Relativity frames), ensuring that passengers stumbling inside the bus, buses driving around the station, and even passengers wandering the station floor cannot exceed this universal speed (Lorentz invariance). Although their motion is random, the passengers and buses still tend to retain a memory of their previous direction (momentum in physics), which means our random walks have a certain persistence. Directional signs posted throughout the station encourage movement towards the center: passengers attempt to cluster near the bus’s center (quantum frame), while the buses, in turn, move towards the station’s center (station’s frame).

The fullness of a bus determines its color. Emptier buses correspond to lower-intensity, lower-frequency red light, while more crowded buses correspond to higher-intensity, higher-frequency blue light. The temperature of the station is related to the total number of passengers present. Just as a hotter oven radiates more blue light, a busier station tends to produce a greater proportion of fully loaded, bluer buses.

In summary:
\begin{itemize}
    \item \textbf{Bus Station:} The blackbody cavity from which energy (passengers) cannot easily escape until it finds its way onto a bus (quantum of light).
    \item \textbf{Buses:} Quanta of energy (like photons), each able to hold a certain number of energy units.
    \item \textbf{Passengers:} Individual energy units that behave like the smallest constituents of light, moving randomly.
    \item \textbf{Drunken Stumbling:} Stochastic random walk (Brownian motion). Both passengers and buses move randomly.
    \item \textbf{Speed Limit:} A universal limit on velocity (akin to the speed of light), ensuring a consistent relativistic frame for all movements.
    \item \textbf{Move-To-Center Signs:} An attractive force guiding energy units to cluster.
    \item \textbf{Bus Fullness and Frequency:} Spectral intensity. Fuller buses emit higher-frequency (bluer) light.
    \item \textbf{Temperature:} The total number of passengers. More passengers correspond to higher temperature, analogous to a hotter blackbody radiating more intense, bluer light.
\end{itemize}


\subsection*{Quantum Gravity, Balls Into Boxes Combinatorics, \& Chaos Strange Attractors}
\addcontentsline{toc}{subsection}{Derivation of Planck's Law \& Strange Attractors}

The key idea is that when energy is quantized into discrete units (like "balls") and placed into modes or states (like "boxes"), the distribution of these units can be understood purely combinatorially. When combined with an attractive force at the quantum scale—analogous to quantum gravity—this leads to stable configurations that yield the well-known Planck distribution without invoking superposition or other paradoxical notions.

\subsubsection*{From Chaos to Stable Quanta: Strange Attractors as Boxes}

Consider a blackbody cavity analogous to a "bus station" filled with many energy quanta ("buses") and discrete energy units ("passengers"). The energy units move randomly (Brownian motion) under a subtle attractive force—akin to a scale-invariant gravity at the quantum level. This force ensures that the quanta form stable "boxes" that can hold a certain number of energy units. In a chaotic, multi-body gravitational system, these stable configurations are like "strange attractors," fractal-like steady states that arise from deterministic chaos.

The presence of an attractive force is crucial: it prevents energy units from dispersing arbitrarily and ensures that stable states (modes) form. In essence, energy units cluster into quanta due to this attractive force, and we can count the probability distributions of energy among these modes using basic combinatorics. This approach turns what is an intractable non-linear dynamics continuous integral problem into an approachable, discrete counting problem.

\subsubsection*{Combinatorial Counting: Balls into Boxes}

We start by modeling the distribution of $M$ indistinguishable energy units (balls) among $N$ distinguishable modes (boxes). The number of ways to do this is given by the combinatorial formula:
\[
W = \binom{N+M-1}{M} = \frac{(N+M-1)!}{M!(N-1)!}.
\]
This formula counts the number of microstates (configurations) of energy distribution.

\subsubsection*{Increasing Temperature: Entropy and Stirling’s Approximation}
In our bus station analogy, increasing temperature corresponds to adding more passengers. Our initial task was to formulate the relationship between the color of departing buses and the total number of passengers in the station. Another way to phrase this is to ask how passengers will distribute themselves among the buses as the station gets busier. In thermodynamics which deals with atoms (much, much larger than quanta) there is a concept of entropy which relates how energy gets distributed. Why entropy should apply at the quantum scale isn't obvious unless quanta (our buses) are made up of smaller energy units (our passengers), however, since that is our assumption it makes sense to use the entropy relationship now. In thermodynamics our buses would be called "microstates" and the different passenger counts of buses are called "modes", so we will use that terminology interchangeably.
\newline

The entropy $S$ is related to the number of microstates by Boltzmann’s relation:
\[
S = k_B \ln(W).
\]
\newline
When $N$ and $M$ are large, we apply Stirling’s approximation:
\[
\ln(n!) \approx n \ln(n) - n.
\]
Using this approximation on $W$, we find a simpler expression for $S$:
\[
S \approx k_B \left[ (N+M)\ln(N+M) - N\ln(N) - M\ln(M)\right].
\]
\newline
Interpreting this entropy combinatorially reveals that the most probable distribution of energy among modes emerges naturally from counting arguments. The complicated integrals from the classical derivation of Planck’s Law are replaced here by discrete, combinatorial reasoning.

\subsubsection*{Relating Entropy to Temperature}
Returning now to the question of "what happens to the color of the buses as the station gets busier," we can ask this question mathematically by looking at the impact of adding a small number of passengers to the station  (i.e. we will take the derivative). Thermodynamics links temperature $T$ to entropy $S$ via:
\[
\frac{1}{T} = \frac{dS}{dU}
\]
where $U$ is the total energy. Differentiating the combinatorial form of $S$ with respect to $U$ and solving for $U$ as a function of $T$ yields a Bose-Einstein-like distribution. Identifying the quantum of energy in each mode as $hf$ (where $h$ is Planck’s constant and $f$ is frequency), we arrive at:
\[
\frac{U}{N} = \frac{hf}{e^{\frac{hf}{k_B T}} - 1}.
\]
\newline
This is the average energy per mode at frequency $f$ and temperature $T$, matching the quantum-statistical form known from standard quantum theory—but derived here from combinatorial logic and an underlying attractive force.

\subsubsection*{Density of States and Planck’s Law}

To obtain the full Planck’s Law, we multiply the average mode energy by the density of states $g(f)$. In a three-dimensional cavity:
\[
g(f) = \frac{8\pi f^2}{c^3}.
\]
\newline
Thus, the spectral energy density is:
\[
\rho(f,T) = g(f)\frac{U}{N} = \frac{8\pi h f^3}{c^3} \frac{1}{e^{\frac{hf}{k_B T}} - 1}.
\]
\newline
This is the celebrated Planck’s Law, now derived without invoking superposition or local hidden variables. Instead, we used:
\begin{itemize}[nosep]
    \item \textbf{Discrete Quanta:} Energy units placed in modes.
    \item \textbf{Combinatorics:} Counting microstates (balls into boxes).
    \item \textbf{Stirling’s Approximation:} Simplifying the factorial expressions.
    \item \textbf{Chaos \& Attraction (Quantum Gravity):} Ensuring stable, discrete modes form from chaotic dynamics.
    \item \textbf{Scale Invariance:} The same logic holds at all scales, from the quantum to the cosmological.
\end{itemize}
\subsubsection*{Connecting Back to Quantum Gravity and Chaos}

The crucial insight is that introducing a scale-invariant attractive force (i.e., a form of quantum gravity) allows for stable, discrete quantum states. The N-Body chaos becomes ordered into "strange attractors," each attractor acting like a stable "box" that can hold integer multiples of $hf$ energy units. The complex dynamical system is thus reduced to a combinatorial counting problem, revealing the Planck distribution naturally.

In essence:
\begin{itemize}[nosep]
    \item Chaos ensures complexity without hidden variables.
    \item Attraction ensures stable states.
    \item Combinatorics and scale invariance do the mathematical heavy lifting.
\end{itemize}

This viewpoint not only resolves paradoxes like EPR without the need for superposition but also shows that Einstein’s intuition about Quantum Mechanics being incomplete was correct. By incorporating quantum gravity and chaos, the full picture emerges, and the strange attractors at the quantum scale explain quantum statistics straightforwardly.


\subsection*{Why Attraction is Essential: Creating Stable Quanta}

Without the attractive force, energy units would never settle into discrete quanta. They would remain scattered, and the problem would remain intractably complex, akin to trying to solve an N-Body gravitational problem with no simplifying assumptions. The presence of quantum gravity at small scales provides a universal attractor, ensuring energy units cluster into stable configurations that can be enumerated. From this counting emerges the Bose-Einstein-like statistics that underlie Planck’s distribution.

In essence, \emph{There is Nothing Without Attraction}. This force molds the chaotic landscape of random walkers into stable boxes whose occupancies we can count. The result is a derivation that is as combinatorially straightforward as distributing people onto buses.

\subsection*{Scale Invariance and the Emergence of Planck’s Law}

Because this approach is scale-invariant, the same combinatorial reasoning applies at all levels. Zooming out, the differences between quantum and classical scales vanish in the logic of counting. At the smallest scales, quantum gravity ensures that energy units settle into stable attractors, just as at larger scales gravitational wells organize matter into stars and galaxies. This self-similarity of nature across scales—fractal in spirit—allows us to derive fundamental quantum laws from a discrete, combinatorial viewpoint.

When we follow the combinatorial steps in detail (see addendum), we find the well-known Planck’s Law formula. No continuous integrals or calculus-based derivations are strictly necessary, only careful counting of states and the recognition that a scale-invariant attractive force provides the stable "boxes" required to define those states. Thus, the celebrated formula:
\[
\rho(f,T) = \frac{8 \pi h f^3}{c^3} \frac{1}{e^{\frac{hf}{k_B T}} - 1}
\]
emerges from a logic chain grounded in chaos, scale invariance, quantum gravity-induced attractors, and combinatorial counting.

\subsection*{Full Planck's Law Under Quantum Gravity Combinatorics Derivation : Detailed Steps}

For those interested in the full derivation steps—from starting with raw combinatorial expressions (balls into boxes) to applying Stirling’s approximation and connecting the discrete results to the continuous limits known in classical physics—the addendum provides a thorough walk-through. At each step, the scale-invariant attractive force and the resulting stable attractors (our "boxes") play a crucial role, transforming a seemingly intractable N-Body gravitational problem into a manageable combinatorial puzzle. This detailed combinatorial path to Planck’s Law demonstrates that, with the right conceptual tools, Einstein’s original quantum insights can be integrated coherently with gravity and Chaos Theory, eliminating the need for paradoxical assumptions in Quantum Mechanics.

\subsubsection*{Setup and Notation}

\begin{itemize}
    \item We consider $N$ identical quanta (or "boxes") into which we distribute $M$ indistinguishable energy units (or "balls"). Each quantum can be viewed as a discrete Relativity frame containing energy units under quantum gravity.
    \item The energy units move randomly, akin to Brownian motion. For small $N$ and $M$, the distribution of energy units per quantum would be approximately binomial. As $N$ and $M$ become very large, this distribution smooths out, effectively converging to a Poisson distribution.
    \item The total energy is $U$. Defining $\mu$ as the average integer number of units per quantum ($M=U/\mu$), we note that at large $N$ and $M$, the variance of the distribution becomes negligible. In this limit, the mode and mean are nearly identical, justifying the integer assumption for $\mu$.
    \item With both $N$ and $M$ large, Stirling’s approximation becomes applicable and simplifies the combinatorial analysis.
\end{itemize}


\subsubsection*{Counting the Number of Microstates $W$}

We need to count the number of ways to distribute $M$ indistinguishable units into $N$ distinguishable boxes (quanta). The standard combinatorial formula for this distribution is:
\[
W = \binom{N + M - 1}{M} = \frac{(N + M - 1)!}{M! (N-1)!}.
\]
\newline
This formula is often derived from a “stars and bars” argument:\newline
- We have $M$ identical items (stars) to be placed into $N$ distinct boxes, separated by $N-1$ bars. \newline
- In total, we have $M$ stars and $N-1$ bars. \newline
- Thus, we arrange $N + M - 1$ objects of which $M$ are identical of one kind and $N-1$ are identical of another kind, giving the above combination formula.\newline

\subsubsection*{Entropy: $S = k_B \ln(W)$}

The entropy is given by the Boltzmann relation:
\[
S = k_B \ln(W).
\]

Substitute $W$:
\[
S = k_B \ln\left(\frac{(N + M - 1)!}{M! (N - 1)!}\right).
\]

We can rewrite this as:
\[
S = k_B \left[ \ln((N + M - 1)!) - \ln(M!) - \ln((N-1)!) \right].
\]

Since $N$ and $M$ are large, we apply Stirling’s approximation:
\[
\ln(n!) \approx n \ln(n) - n.
\]

Applying Stirling’s approximation to each factorial:

\[
\ln((N+M-1)!) \approx (N + M - 1)\ln(N+M-1) - (N + M - 1),
\]
\[
\ln(M!) \approx M \ln(M) - M,
\]
\[
\ln((N-1)!) \approx (N-1)\ln(N-1) - (N-1).
\]

For large $N$ and $M$, the “-1” terms are negligible compared to $N$ and $M$, so we approximate:
\[
N + M - 1 \approx N + M, \quad N - 1 \approx N.
\]

Thus:
\[
\ln((N + M - 1)!) \approx (N + M) \ln(N + M) - (N + M),
\]
\[
\ln((N-1)!) \approx N \ln(N) - N.
\]

Substitute these approximations back into the entropy expression:
\[
S \approx k_B \left[ (N + M)\ln(N + M) - (N + M) - (M \ln(M) - M) - (N \ln(N) - N) \right].
\]

Combine like terms where tje linear terms $-(N+M) + M + N$ cancel out. So we have a cleaner expression:
\[
S \approx k_B \left[ (N + M)\ln(N + M) - M \ln(M) - N \ln(N) \right].
\]


Recall that $M = U/\mu$, substitute for $M$:
\[
S \approx k_B \left[ (N + \tfrac{U}{\mu})\ln\left(N + \tfrac{U}{\mu}\right) - N \ln(N) - \tfrac{U}{\mu}\ln\left(\tfrac{U}{\mu}\right) \right].
\]

\subsection*{Relating Temperature and Entropy: $\frac{1}{T} = \frac{dS}{dU}$}

Temperature $T$ is defined via the thermodynamic relation:
\[
\frac{1}{T} = \frac{dS}{dU}.
\]

To find $\frac{dS}{dU}$, we treat $N$ and $\mu$ as constants. Differentiate $S$ with respect to $U$:

\[
S \approx k_B \left[ (N + \tfrac{U}{\mu})\ln\left(N + \tfrac{U}{\mu}\right) - N \ln(N) - \tfrac{U}{\mu}\ln\left(\tfrac{U}{\mu}\right)\right].
\]

Focus on the $U$-dependent terms:
\[
\frac{dS}{dU} = k_B \left[ \frac{d}{dU}\left((N + \tfrac{U}{\mu})\ln\left(N + \tfrac{U}{\mu}\right)\right) - \frac{d}{dU}\left(\tfrac{U}{\mu}\ln\left(\tfrac{U}{\mu}\right)\right) \right].
\]

Compute each derivative:\newline
1. For $(N + \tfrac{U}{\mu})\ln\left(N + \tfrac{U}{\mu}\right)$:
\[
\frac{d}{dU} \left((N + \tfrac{U}{\mu})\ln\left(N + \tfrac{U}{\mu}\right)\right) 
= \ln\left(N + \tfrac{U}{\mu}\right)\frac{d}{dU}(N + \tfrac{U}{\mu}) + (N+\tfrac{U}{\mu}) \frac{1}{N+\tfrac{U}{\mu}}\frac{d}{dU}(N+\tfrac{U}{\mu})
\]
\[
= \ln\left(N + \tfrac{U}{\mu}\right)\frac{1}{\mu} + \frac{N+\tfrac{U}{\mu}}{N+\tfrac{U}{\mu}}\frac{1}{\mu}
= \frac{1}{\mu}\ln\left(N + \tfrac{U}{\mu}\right) + \frac{1}{\mu}.
\]
\newline
2. For $\tfrac{U}{\mu}\ln\left(\tfrac{U}{\mu}\right)$:
\[
\frac{d}{dU}\left(\tfrac{U}{\mu}\ln\left(\tfrac{U}{\mu}\right)\right) 
= \ln\left(\tfrac{U}{\mu}\right)\frac{1}{\mu} + \tfrac{U}{\mu}\frac{1}{U}\frac{1}{\mu}U\biggr| \text{(chain rule details)}.
\]

More carefully:
\[
\frac{d}{dU}\left(\frac{U}{\mu}\ln\left(\frac{U}{\mu}\right)\right) 
= \frac{1}{\mu}\ln\left(\frac{U}{\mu}\right) + \frac{U}{\mu}\frac{1}{U}
= \frac{1}{\mu}\ln\left(\frac{U}{\mu}\right) + \frac{1}{\mu}.
\]

Putting it all together:
\[
\frac{dS}{dU} = k_B\left[\left(\frac{1}{\mu}\ln\left(N + \frac{U}{\mu}\right) + \frac{1}{\mu}\right) - \left(\frac{1}{\mu}\ln\left(\frac{U}{\mu}\right) + \frac{1}{\mu}\right)\right].
\]

The $+\frac{1}{\mu}$ and $-\frac{1}{\mu}$ cancel out:
\[
\frac{dS}{dU} = \frac{k_B}{\mu}\left[\ln\left(N + \frac{U}{\mu}\right) - \ln\left(\frac{U}{\mu}\right)\right].
\]

Use the logarithm property $\ln(a) - \ln(b) = \ln\left(\frac{a}{b}\right)$:
\[
\frac{dS}{dU} = \frac{k_B}{\mu}\ln\left(\frac{N + \frac{U}{\mu}}{\frac{U}{\mu}}\right) = \frac{k_B}{\mu}\ln\left(1 + \frac{N\mu}{U}\right).
\]

We have:
\[
\frac{1}{T} = \frac{dS}{dU} = \frac{k_B}{\mu}\ln\left(1 + \frac{N\mu}{U}\right).
\]

\subsection*{Solving for $U$ in Terms of $T$}

Rewrite:
\[
\frac{\mu}{k_B T} = \ln\left(1 + \frac{N\mu}{U}\right).
\]

Exponentiate both sides:
\[
e^{\frac{\mu}{k_B T}} = 1 + \frac{N\mu}{U}.
\]

Subtract 1:
\[
e^{\frac{\mu}{k_B T}} - 1 = \frac{N\mu}{U}.
\]

Invert:
\[
U = \frac{N\mu}{e^{\frac{\mu}{k_B T}} - 1}.
\]

\subsubsection*{Average Energy per Quantum (Mode)}

The above $U$ is the total energy. The average energy per quantum (or mode) is:
\[
\frac{U}{N} = \frac{\mu}{e^{\frac{\mu}{k_B T}} - 1}.
\]

Recall $\mu$ is the mode number of energy units within a quantum associated with that frequency. For electromagnetic radiation, Planck identified this quantum of energy as:
\[
\mu = hf,
\]
where $h$ is Planck’s constant and $f$ is the frequency.

Substitute $\mu = hf$:
\[
\frac{U}{N} = \frac{hf}{e^{\frac{hf}{k_B T}} - 1}.
\]

This is the mean energy of an oscillator (mode) at frequency $f$ and temperature $T$.

\subsection*{Deriving the Full Planck’s Law for Spectral Energy Density}

Planck’s Law for the spectral energy density $\rho(f,T)$ is related to the average energy per mode by the density of states. In a 3D cavity of volume $V$, the number of modes per unit volume per unit frequency is given by:
\[
g(f) = \frac{8\pi f^2}{c^3}.
\]

Multiply the average energy per mode by $g(f)$ to get the energy density:
\[
\rho(f,T) = g(f)\frac{U}{N} = \frac{8\pi f^2}{c^3} \cdot \frac{hf}{e^{\frac{hf}{k_B T}} - 1}.
\]

Combine factors:
\[
\rho(f,T) = \frac{8\pi h f^3}{c^3}\frac{1}{e^{\frac{hf}{k_B T}} - 1}.
\]

This is the standard form of Planck’s Law for the spectral energy density of blackbody radiation.


\subsection*{Summary of the Derivation Steps}

\begin{enumerate}
    \item Start with the combinatorial expression for the number of ways $W$ to distribute energy units.
    \item Express entropy $S = k_B\ln(W)$.
    \item Use Stirling’s approximation to find a closed form for $S$ in terms of $U$.
    \item Differentiate $S$ w.r.t. $U$ to find $\frac{1}{T}$.
    \item Solve the resulting equation for $U$, giving the Bose-Einstein-like distribution formula $U = \frac{N\mu}{e^{\mu/(k_B T)}-1}$.
    \item Identify $\mu = hf$ for photons, yielding the average energy per mode.
    \item Multiply by the known density of states $g(f)$ to get Planck’s Law.
\end{enumerate}

This detailed derivation completes the purely discrete, combinatorial, and scale-invariant approach to obtaining Planck’s Law under quantum gravity.




\section*{There is Nothing Without Attraction: Scale-Invariant Attractive Force}
\addcontentsline{toc}{section}{There is Nothing Without Attraction: Scale-Invariant Attractive Force}

In order to bring gravity into Quantum Mechanics and possibly complete the theory we had two approaches to choose from-calculus or discrete mathematics. Recalling that the N-Body problem, involving three or more bodies attracted to one another, is famously resistant to calculus solutions, we ruled calculus out as a mathematical form that Einstein would have sought. Instead, we found that introducing a scale-invariant attractive force gives us stable configurations of quanta so it transforms the problem into one of counting—effectively turning the complex dynamics into a balls-into-boxes scenario. Attraction ensures that the energy units form coherent quanta (our “boxes”), making a combinatorial approach feasible.

Without attraction to balance the chaos inducing entropy, energy units would never coalesce into discrete quanta, and we would face an intractable problem. With attraction, order emerges from chaos: the presence of an attractive force balances Brownian motion, ensuring that while every configuration is possible, the most likely configurations are those clustered around stable centers. In the bus station analogy, it is the attraction that gives us well-defined “buses” and “station center,” enabling us to count distributions and ultimately derive the laws governing the emission spectrum.

\subsection*{Unifying the Fundamental Forces Through Scale Invariance}
\addcontentsline{toc}{subsection}{Unifying the Fundamental Forces Through Scale Invariance}

Having established a mathematical bridge between Quantum Mechanics and thermodynamics across an enormous scale range (on the order of $10^{25}$ in magnitude), it is natural to ask whether similar connections hold at even larger scales. The Anti-de Sitter/Conformal Field Theory (AdS/CFT) correspondence provides a well-known theoretical framework in which scale invariance is central. AdS/CFT relates a gravity theory within a bulk spacetime (often containing black holes) to a lower-dimensional, scale-invariant quantum field theory defined on its boundary. A key holographic insight from AdS/CFT is that the internal content of a black hole is fully described by the entropy on its boundary surface. Intriguingly, this is conceptually similar to our discrete combinatorial approach, where the content of a quantum is likewise reflected in the entropy on its “surface.”

While the conventional treatments of AdS/CFT rely heavily on continuous mathematics and advanced field theory techniques, the essential entropy relationship it reveals can also be understood as a discrete, scale-invariant correlation of the type we have just derived. In other words, AdS/CFT demonstrates that General Relativity—at the scale of black holes—is also governed by principles of scale invariance. This parallels the logic used to unify Quantum Mechanics and thermodynamics, further strengthening the argument that all of the fundamental forces may emerge from a single, scale-invariant framework.

This perspective suggests that the commonly accepted hierarchy of fundamental interactions (strong, weak, electromagnetic, and gravitational) might be interpreted as scale-dependent manifestations of one underlying attractive principle. Quantum gravity and Modified Newtonian Dynamics (MOND) can be included as bookends of this hierarchy. Thus, what appear to be distinct physical theories at different scales are likely different frames of the same combinatorial, scale-invariant structure.

\subsubsection*{Perceived Attraction Strength and Frame Hierarchy}
\addcontentsline{toc}{subsubsection}{Perceived Strength and Frame Hierarchy}

Consider a nested hierarchy of frames:
\[
F_0 \subset F_1 \subset F_2 \subset \cdots \subset F_M,
\]
where $F_N$ denotes a frame at a given dimensional or scale level $N$, and $M > N$ indicates a higher-level frame encompassing a larger structure. Each frame contains smaller sub-frames, much as atoms nest within molecules, which nest within cells, and so forth, eventually reaching cosmic scales.

Let $b$ be a body within a frame $F_N^{(i)}$ at scale $N$. Internally, $b$ experiences a deterministic, discrete "pull" of its constituent parts toward the center of $F_N^{(i)}$, representing the fundamental attractive action at that scale. An observer residing at a higher scale $M > N$ does not directly resolve these microscopic details. Instead, they perceive only an aggregated effect: the numerous strong, small-scale attractions average out to produce what, at their coarser resolution, appears as a weaker, more diffuse force.

In short, as the observational scale enlarges, the same fundamental attraction—when averaged over vast ensembles—appears weaker and more long-ranged. At the smallest scales, where $L_N$ (the characteristic length scale at level $N$) is tiny, the attraction manifests as a very strong force. At larger scales, where $L_N$ is enormous, the same principle yields a weaker effective interaction, eventually aligning with our observation that gravity is the weakest known force.

\subsubsection*{A Loose Mathematical Formalization}
\addcontentsline{toc}{subsubsection}{A Loose Mathematical Formalization}

Define $\Delta r_N$ as the fixed, discrete shift magnitude (e.g., a small number of “pixels” in the discrete model) that particles within frame $F_N$ move toward the center in each step. Let $L_N$ be the characteristic length scale associated with frame $F_N$. We can then define an effective “attraction measure” at scale $N$ as:
\[
A_N = \frac{\Delta r_N}{L_N}.
\]

Since $\Delta r_N$ is fixed, increasing $L_N$ reduces $A_N$. Thus,
\[
A_N \propto \frac{1}{L_N}.
\]

At very small $L_N$, $A_N$ is large, reflecting the strong forces we see at subatomic scales. As $L_N$ grows (moving to atomic, molecular, macroscopic, and ultimately cosmological scales), $A_N$ diminishes, mirroring the observed weakening of forces like gravity.

\subsubsection*{Comparison to Known Forces}
\addcontentsline{toc}{subsubsection}{Comparison to Known Forces}

\begin{itemize}
    \item \textbf{Quantum Gravity:}  
    At the smallest conceivable scales, $L_N$ is extremely small, making $A_N$ effectively enormous. The constituents appear point-like and indivisible, consistent with the idea of fundamental quanta of gravity.

    \item \textbf{Strong Force:}  
    Acts at very small scales, where $L_N$ is minimal and $A_N$ is large, matching the strong binding energy inside atomic nuclei.

    \item \textbf{Weak Force:}  
    At slightly larger scales than the strong force, $L_N$ is bigger, so $A_N$ decreases, aligning with the relatively weaker but still short-range character of the weak interaction.

    \item \textbf{Electromagnetic Force:}  
    Operating at atomic and molecular scales, $L_N$ is larger still. $A_N$ is therefore smaller than at nuclear scales, though still stronger than gravity, which emerges at even larger $L_N$.

    \item \textbf{Gravitational Force:}  
    Dominant at macroscopic to cosmological scales, where $L_N$ is vast. Thus, $A_N$ is extremely small, corresponding to the gravitational force’s relative weakness.

    \item \textbf{MOND:}  
    At the largest galactic and intergalactic scales, $L_N$ becomes so large that subtle modifications to Newtonian dynamics are observed (MOND). These are naturally explained as minor adjustments to $A_N$ under extraordinary scale expansion.
\end{itemize}

In this picture, all forces arise from the same underlying scale-invariant principle, with their differences in strength and range a direct consequence of the frame hierarchy and scale. By applying the same combinatorial, discrete logic used to unify Quantum Mechanics and thermodynamics, one obtains a coherent, scale-invariant narrative spanning from quantum gravity to galactic dynamics.

\section*{Conclusion And Computational Implications}
\addcontentsline{toc}{section}{Conclusion}

By integrating the lessons of the blackbody problem and the thermodynamic-like combinatorial approach with a scale-invariant attractive force, we produce a coherent framework in which the observed hierarchy of fundamental forces becomes a natural consequence of nested frame structures. The blackbody problem, initially a puzzle of radiation and energy distribution, revealed that discrete combinatorics and attraction suffice to derive Planck’s Law. Scaling this idea upwards and downwards across the frame hierarchy explains how the same fundamental rule of attraction manifests at different scales with radically different "strengths."

In this view, quantum gravity is simply another name for the same fundamental force acting at the smallest scales, while gravity at large scales is that force averaged and attenuated through countless layers of frame nesting. The thermodynamic ensembles that appear at one scale and the Boltzmann-like distributions at another scale both arise from the same discrete, scale-invariant principles. Thus, scale invariance offers an elegant, unifying perspective, potentially guiding future research in quantifying these scaling relations precisely and exploring testable predictions across domains from particle physics to cosmology.

We started with the notion that the most valuable computations are physics computations of highly complex systems of the type superposition based quantum computing is supposed to solve. In showing that Einstein's EPR paradox can be resolved by a more complete understanding of physics that incorporates quantum gravity, it becomes highly unlikely that superposition based quantum computers will work on the principles they are supposed to. However, with a more complete understanding of physics, and an eye to the compression implications of self-similar fractal structures, we can begin to see how a new and highly effecient computational approach might be possible.
\newpage

\bibliographystyle{plain}
\bibliography{references}

\end{document}
